\documentclass[11pt]{article}
\usepackage[T1]{fontenc}
\usepackage[utf8]{inputenc}
\usepackage{lmodern}
\usepackage[a4paper,margin=28mm]{geometry}
\usepackage{amsmath,amssymb,amsthm,mathtools}
\usepackage{enumitem,booktabs,array,microtype}
\usepackage{tikz-cd}
\usepackage{xcolor}
\usepackage[unicode,pdfencoding=auto,psdextra,hyperfootnotes=false]{hyperref}
\usepackage[nameinlink,capitalise,noabbrev]{cleveref}
\hypersetup{colorlinks=true,linkcolor=blue,citecolor=blue,
 urlcolor=blue,
 pdftitle={A note of refinements to an algebra of infinite little finitistic dimension},
 pdfauthor={Qiyue Tang},pdfsubject={A finite-dimensional algebra of infinite little finitistic dimension}}
\urlstyle{same}
\pdfgentounicode=1
\pdfinfoomitdate=1
\pdftrailerid{}
\pdfsuppressptexinfo=15
\newtheorem{theorem}{Theorem}[section]
\numberwithin{equation}{section}
\newtheorem{proposition}[theorem]{Proposition}
\newtheorem{lemma}[theorem]{Lemma}
\newtheorem{corollary}[theorem]{Corollary}
\theoremstyle{definition}
\newtheorem{definition}[theorem]{Definition}
\theoremstyle{remark}
\newtheorem{remark}[theorem]{Remark}
\AddToHook{env/theorem/begin}{\crefalias{theorem}{theorem}}
\AddToHook{env/proposition/begin}{\crefalias{theorem}{proposition}}
\AddToHook{env/lemma/begin}{\crefalias{theorem}{lemma}}
\AddToHook{env/corollary/begin}{\crefalias{theorem}{corollary}}
\AddToHook{env/definition/begin}{\crefalias{theorem}{definition}}
\AddToHook{env/remark/begin}{\crefalias{theorem}{remark}}
\crefname{theorem}{Theorem}{Theorems}
\Crefname{theorem}{Theorem}{Theorems}
\crefname{proposition}{Proposition}{Propositions}
\Crefname{proposition}{Proposition}{Propositions}
\crefname{lemma}{Lemma}{Lemmas}
\Crefname{lemma}{Lemma}{Lemmas}
\crefname{corollary}{Corollary}{Corollaries}
\Crefname{corollary}{Corollary}{Corollaries}
\newcommand{\C}{\mathbb C}
\newcommand{\F}{\mathbb F}
\newcommand{\Z}{\mathbb Z}
\newcommand{\N}{\mathbb N}
\newcommand{\D}{\mathsf D}
\newcommand{\K}{\mathsf K}
\newcommand{\otL}{\otimes^{\mathbf L}}
\DeclareMathOperator{\Hom}{Hom}
\DeclareMathOperator{\End}{End}
\DeclareMathOperator{\Ext}{Ext}
\DeclareMathOperator{\Tor}{Tor}
\DeclareMathOperator{\pd}{pd}
\DeclareMathOperator{\gldim}{gl.dim}
\DeclareMathOperator{\findim}{fin.dim}
\DeclareMathOperator{\Findim}{Fin.dim}
\DeclareMathOperator{\rad}{rad}
\DeclareMathOperator{\id}{id}
\DeclareMathOperator{\Cone}{Cone}
\DeclareMathOperator{\Kar}{Kar}
\DeclareMathOperator{\supp}{supp}
\DeclareMathOperator{\thick}{thick}
\DeclareMathOperator{\perf}{perf}
\DeclareMathOperator{\proj}{proj}
\DeclareMathOperator{\Mod}{Mod}
\DeclareMathOperator{\modu}{mod}
\setlist[enumerate]{label=\textup{(\roman*)},leftmargin=2em}
\setlength{\emergencystretch}{1.5em}
\allowdisplaybreaks[1]

\usepackage{indentfirst}
\setlength{\parindent}{1.5em}

\title{A note of refinements to an algebra of infinite little finitistic dimension}
\author{Qiyue Tang}
\date{\small Generated by GPT-6-Astra\\[6pt]
October 7, 2026\\[3pt]\small This paper is not intended for journal submission.}
\begin{document}
\maketitle
\begingroup\renewcommand{\thefootnote}{}\footnotetext{\textit{2020 Mathematics Subject Classification.}
Primary 16E10; Secondary 16G20.
\textit{Keywords.} Finitistic dimension, trivial extension, derived tensor product, brick.}\endgroup
\begin{abstract}
Building on OpenAI's \emph{An algebra of infinite little finitistic dimension},
we refine the construction and extend it to every field.
\end{abstract}
\section{Introduction}\label{sec:introduction}
For a finite-dimensional algebra $A$ over a field $k$, write $\pd_A N$
for the projective dimension of a left $A$-module $N$.
Throughout, the projective dimension of the zero module or zero complex
is $-\infty$, the supremum of the empty set is $-\infty$, and
$-\infty+n=-\infty$ for every integer $n$. Set
\[
\findim A=\sup\{\pd_A N:N\text{ is a finite-dimensional left }A\text{-module},\ \pd_A N<\infty\}.
\]
The finitistic-dimension problem asks whether this supremum is finite
for every finite-dimensional algebra; see \cite{Bass1960,HuisgenZimmermann1995}.
OpenAI's construction \cite{OpenAI198} gives a complex algebra for which
the supremum is infinite. We retain its group, localization, tensor
realization, and square-zero extension constructions, and refine their
field of definition, test modules, and dimension estimates.

For an algebra $A$, let $J_A=\rad A$ denote its Jacobson radical.
The Loewy length of a nonzero module $N$ is the least $h\ge1$ with
$J_A^hN=0$. An algebra is split basic if $A/J_A$ is a finite product
of copies of $k$. We call a module a brick when its endomorphism ring
is a division ring; the complex bricks below have endomorphism ring $\C$.

\begin{theorem}\label{thm:main}
For every field $k$ there are a finite-dimensional split basic
$k$-algebra $A$ with $J_A^7=0$, integers $b_*\ge0$ and $l\ge0$,
and finite-dimensional left $A$-modules $N_m$, $m\ge1$, such that
\[
(b_*+5)(m-1)+1\le\pd_A N_m\le(b_*+5)(m-1)+l+3.
\]
The modules have Loewy length three, their endomorphism rings are
division rings, and $\Hom_A(N_m,N_n)=0$ for $m\ne n$.
In particular, $\findim A=\infty$.
For $k=\C$ they can be chosen with
\[
\dim_{\C}N_m=3\cdot2^{3m},\qquad \End_A(N_m)=\C.
\]
Every primitive vertex idempotent $\varepsilon$ in this construction
satisfies $\varepsilon A\varepsilon=k\varepsilon$.
\end{theorem}

The proof follows the original construction in \cite{OpenAI198}.
A finitely presented group supplies a group algebra $R$, a central
idempotent $e$, and an automorphism $\alpha$. For a left $R$-module
$Y$, the operation $H(Y)={}_{\alpha}(eY)$ takes the summand $eY$ and lets
$r\in R$ act by $\alpha(r)$. We construct simple modules that vanish
under precisely the $m$th iteration of $H$.
A three-vertex algebra $B$ encodes $R$, and a bounded bimodule complex
$P$ realizes $H$ together with a shift by three.
A chain algebra records the differential of $P$ in an ordinary
bimodule $X$ over an algebra $D$ of finite global dimension.
The square-zero extension of $D$ by $X$ converts finite tensor
extinction into finite projective dimension.

All algebras are associative and unital, and algebra homomorphisms
preserve $k$ and the unit. Modules are left modules unless designated
right. The notation $\D^b(A\text{-}\modu)$ means the bounded derived
category of finite-dimensional left $A$-modules; $\K^b(\proj\text{-}A)$
means the homotopy category of bounded finitely generated projective
right $A$-module complexes. The symbol $\otL$ denotes derived tensor
product. For a complex $K$ and an integer $n$, write
$d_K^n:K^n\to K^{n+1}$ for its
differential and $H^n(K)=\ker d_K^n/\operatorname{im}d_K^{n-1}$
for its cohomology. For integers $n,s$, the shift convention is
\[
K[s]^n=K^{n+s},\qquad d_{K[s]}^n=(-1)^s d_K^{n+s}.
\]
A module is regarded as a complex concentrated in degree zero, so
its shift by $s$ has cohomology in degree $-s$.
Path products and endomorphism products are composed from right to left.

\section{A selection process with unbounded extinction}
\label{sec:selection}

Let $R$ be a unital $k$-algebra, let $e\in R$ be a central idempotent,
and let $\alpha:R\to R$ be a unital $k$-algebra endomorphism.  For a left
$R$-module $Y$, define
\begin{equation}
  H(Y)={}_{\alpha}(eY).
  \label{eq:selection-functor}
\end{equation}
Here $eY$ is an $R$-submodule, and the subscript denotes restriction of
scalars along $\alpha$: on the vector space $eY$, an element $r\in R$
acts by the original operator $\alpha(r)$. On a module map $f:Y\to Z$,
the functor acts by the restriction $f|_{eY}:eY\to eZ$, with both actions
twisted by $\alpha$. It is additive and preserves finite dimensionality.
For $j\ge0$, write $H^{\circ j}$ for its $j$-fold iterate, with
$H^{\circ 0}$ the identity functor; superscripts without a circle on
$H$ will denote cohomology of a complex.

\begin{proposition}
\label{prop:selection}
For every field $k$, there exist a finitely presented unital $k$-algebra $R$, a central
idempotent $e\in R$, and an automorphism $\alpha$ of $R$ such that, for
every integer $m\geq 1$, a finite-dimensional simple left $R$-module $Y_m$
satisfies
\begin{equation}
  H^{\circ j}(Y_m)\neq 0\quad(0\leq j<m),
  \qquad H^{\circ m}(Y_m)=0.
  \label{eq:extinction}
\end{equation}
\end{proposition}

Fix a prime $p$ invertible in $k$: take $p=2$ if $\operatorname{char}k\ne2$,
and $p=3$ otherwise. We construct a finitely presented group $G$
with central elements of order $p$ and take its group algebra $R=kG$.

The group presentation is modeled on matrix groups of Abels type.
The commutator transfer and centrality argument parallel
\cite[Proposition~4.9, equations~(4.9)--(4.10)]{Rego2022}.

\subsection{The group and its finite presentation}

Write $[x,y]=xyx^{-1}y^{-1}$. Define a group $G$ with commuting
generators $T_\ell$, $\ell\in\{1,2,3\}$, and generators
\[
  U_i(r),\qquad V_i(r),\qquad W_{ij}(r)
  \quad(i,j\in\{1,2,3\},\ i\neq j,\ r\in\Z).
\]
Each generator type $S$ has a weight $\lambda_S\in\Z^3$, given by
\begin{equation}
  \lambda_{U_i}=-\varepsilon_i,\qquad
  \lambda_{V_i}=\varepsilon_i,\qquad
  \lambda_{W_{ij}}=\varepsilon_i-\varepsilon_j,
  \label{eq:group-weights}
\end{equation}
where $\varepsilon_1,\varepsilon_2,\varepsilon_3$ is the standard basis
of $\Z^3$. Write $\lambda_{S,\ell}$ for the $\ell$th coordinate
of $\lambda_S$. Impose the conjugation rules
\begin{equation}
  T_\ell S(r)T_\ell^{-1}=S(r+\lambda_{S,\ell}).
  \label{eq:group-conjugation}
\end{equation}
The remaining relations, with all arguments ranging independently
over $\Z$, are
\begin{enumerate}
  \item $U_i(r)^p=1$;
  \item generators of types $U_i,U_j$, of types $V_i,V_j$, and of
    types $U_i,V_j$ commute whenever $i\neq j$;
  \item $W_{ij}$ commutes with $U_h$ for $h\neq i$ and with $V_h$
    for $h\neq j$;
  \item for $i\neq j$,
  \begin{equation}
    [U_i(r),W_{ij}(s)]=U_j(r+s),\qquad
    [W_{ij}(s),V_j(r)]=V_i(s+r).
    \label{eq:group-commutator-relations}
  \end{equation}
\end{enumerate}

\begin{lemma}
\label{lem:group-finite-presentation}
The group $G$ is finitely presented.
\end{lemma}

\begin{proof}
For $n=(n_1,n_2,n_3)\in\Z^3$, put
$T^n=T_1^{n_1}T_2^{n_2}T_3^{n_3}$ and regard a weight as the map
$\lambda_S(n)=\sum_\ell\lambda_{S,\ell}n_\ell$.  Retain only the
fifteen generators
\[
  T_\ell,\quad u_i=U_i(0),\quad v_i=V_i(0),\quad
  w_{ij}=W_{ij}(0)\quad(i\neq j).
\]
Besides commutation of the $T_\ell$, require the following finite
set of centralizer relations.  Each of $u_i,v_i$ commutes with
$T_h$ for $h\neq i$, and $w_{ij}$ commutes with $T_iT_j$ and with
$T_h$ for the unique $h\notin\{i,j\}$.  These are precisely
centralization by a basis of each kernel lattice:
\[
\begin{aligned}
 \ker\lambda_{U_i}=\ker\lambda_{V_i}
   &=\bigoplus_{h\neq i}\Z\varepsilon_h,\\
 \ker\lambda_{W_{ij}}
   &=\Z(\varepsilon_i+\varepsilon_j)
     \oplus\Z\varepsilon_h
       \quad(h\notin\{i,j\}).
\end{aligned}
\]
Since the $T_\ell$ commute, the kernel lattice also centralizes
every $T$-conjugate of its argument-zero generator.  Thus one may define
\begin{equation}
\begin{aligned}
  U_i(r)&=T_i^{-r}u_iT_i^r,&
  V_i(r)&=T_i^rv_iT_i^{-r},&
  W_{ij}(r)&=T_i^rw_{ij}T_i^{-r}.
\end{aligned}
\label{eq:group-finite-generators}
\end{equation}
Equivalently, $S(r)=T^nS(0)T^{-n}$ for any $n$ with
$\lambda_S(n)=r$.  The kernel relations make this independent of
the choice of $n$ and imply all of \eqref{eq:group-conjugation}.

It remains to impose the argument-zero instance of each of the
remaining relation patterns.  For every pair of input types
occurring in those patterns, the two weights give a surjection
$\Z^3\to\Z^2$.  The following table supplies integral preimages of
an arbitrary pair $(r,s)$; unmentioned coordinates are zero.
\begin{center}
\small
\begin{tabular}{@{}lll@{}}
\toprule
Input types & Index condition & Coordinates of $n$ \\
\midrule
$U_i,U_j$ & $i\neq j$ & $n_i=-r,\ n_j=-s$ \\
$V_i,V_j$ & $i\neq j$ & $n_i=r,\ n_j=s$ \\
$U_i,V_j$ & $i\neq j$ & $n_i=-r,\ n_j=s$ \\
$W_{ij},U_j$ & $i\neq j$ & $n_j=-s,\ n_i=r-s$ \\
$W_{ij},U_h$ & $h\notin\{i,j\}$ &
  $n_i=r,\ n_j=0,\ n_h=-s$ \\
$W_{ij},V_i$ & $i\neq j$ & $n_i=s,\ n_j=s-r$ \\
$W_{ij},V_h$ & $h\notin\{i,j\}$ &
  $n_i=r,\ n_j=0,\ n_h=s$ \\
$U_i,W_{ij}$ & $i\neq j$ & $n_i=-r,\ n_j=-r-s$ \\
$W_{ij},V_j$ & $i\neq j$ & $n_j=s,\ n_i=r+s$ \\
\bottomrule
\end{tabular}
\end{center}
Conjugating a zero-argument commutation relation by the corresponding
$T^n$ therefore gives every required instance.  For the two relations
in \eqref{eq:group-commutator-relations}, the output weight is the sum
of the input weights:
\[
  -\varepsilon_i+(\varepsilon_i-\varepsilon_j)=-\varepsilon_j,
  \qquad
  (\varepsilon_i-\varepsilon_j)+\varepsilon_j=\varepsilon_i.
\]
Hence the same conjugation also gives the correct output argument
$r+s$.  The power relation propagates from argument zero because
each individual weight is surjective onto $\Z$.

We have obtained all the defining relations from finitely many
relators on fifteen generators.  Conversely, the chosen kernel
centralizers and zero-argument relators hold in the original
presentation by \eqref{eq:group-conjugation}.  The two presentations
are therefore equivalent.
\end{proof}

\subsection{Central elements and an automorphism}

Central elements shifted by an automorphism will determine the
extinction time.

\begin{lemma}
\label{lem:central-shift}
For each $N\in\Z$, the expression
\begin{equation}
  z_N=[U_i(a),V_i(b)],\qquad a+b=N,
  \label{eq:central-involutions}
\end{equation}
is independent of $i,a,b$.  Each $z_N$ is central and satisfies
$z_N^p=1$.  There is an automorphism $\alpha_G$ of $G$ fixing all
$T_\ell,V_i(r),W_{ij}(r)$ and sending $U_i(r)$ to $U_i(r+1)$;
it satisfies $\alpha_G(z_N)=z_{N+1}$.
\end{lemma}

\begin{proof}
Fix distinct $i,j$ and integers $r,q,s$, and set
\[
 x=U_i(r),\quad y=W_{ij}(q),\quad v=V_j(s),\qquad
 u=[x,y]=U_j(r+q),\quad w=[y,v]=V_i(q+s).
\]
The defining commutations give $[x,v]=[w,u]=[w,v]=1$.
Using $xy=uyx$ and $yv=wvy$, reorder $xyv$ in the two ways
\[
  xyv=uwvyx=(xwx^{-1})vuyx.
\]
Cancellation gives $uwv=[x,w]wvu$, and consequently
\[
  [x,w]=uwvu^{-1}v^{-1}w^{-1}=uvu^{-1}v^{-1}=[u,v].
\]
The last equality uses only the commutation of $w$ with $u$ and $v$.
We have proved
\begin{equation}
 [U_i(r),V_i(q+s)]=[U_j(r+q),V_j(s)].
 \label{eq:commutator-transfer}
\end{equation}
Given $a+b=c+d=N$ at distinct internal indices, choose
$r=a$, $s=d$, and $q=c-a=b-d$ in
\eqref{eq:commutator-transfer}.  This compares the two expressions
in \eqref{eq:central-involutions}.  Two splittings at the same
internal index can be compared through another index, proving
the claimed independence.

To show that $z_N$ commutes with $U_h(r)$ or $V_h(r)$, express
$z_N$ using an index $i\neq h$.  Both factors of that commutator
commute with the tested generator.  To commute with $W_{hj}(r)$,
use the index outside $\{h,j\}$, so that again both factors commute
with it. Write $\delta_{ab}=1$ if $a=b$ and $0$ otherwise.
Conjugation by $T_\ell$ changes the arguments
$(a,b)$ in \eqref{eq:central-involutions} to
$(a-\delta_{\ell i},b+\delta_{\ell i})$.  Their sum is unchanged,
so the splitting independence proves invariance under $T_\ell$.
Thus $z_N$ is central.  Since $U_i(a)^p=1$, centrality also gives
\[
  1=[U_i(a)^p,V_i(b)]=[U_i(a),V_i(b)]^p=z_N^p.
\]

Translation of every $U$-argument by $1$ preserves the conjugation,
power, and commutation relations.  The first relation in
\eqref{eq:group-commutator-relations} becomes
\[
  [U_i(r+1),W_{ij}(s)]=U_j(r+s+1),
\]
which is another defining instance; the second relation is unchanged.
This defines an endomorphism $\alpha_G$.  Translation by $-1$
preserves the relations for the same reason and is its two-sided
inverse.  Formula \eqref{eq:central-involutions} then gives
$\alpha_G(z_N)=z_{N+1}$.
\end{proof}

\subsection{Finite quotients detecting independent central elements}

\begin{lemma}
\label{lem:finite-central-quotients}
For every $m\geq1$, there is a finite quotient $F_m$ of $G$ in which
the images of $z_0,\ldots,z_{m-1}$ generate a central subgroup
isomorphic to $(\Z/p\Z)^m$.
\end{lemma}

\begin{proof}
Let $\F_p$ be the field with $p$ elements and put
$S_m=\F_p[t]/(t^m-1)$, using $t$ also for its residue class.
This ring is finite, and $t$ is a unit
with inverse $t^{m-1}$, so $t^r$ is defined for every $r\in\Z$.
Use $5\times5$ matrices over $S_m$, indexed by $0,1,2,3,4$.
Let $I_5$ be the identity matrix and $E_{ab}$ the matrix with entry
$1$ at $(a,b)$ and $0$ elsewhere. Send $T_\ell$ to the diagonal
matrix with $t$ in position $\ell$ and $1$ elsewhere, and set
\begin{equation}
\begin{aligned}
 U_i(r)&\longmapsto I_5+t^rE_{0i},&
 V_i(r)&\longmapsto I_5+t^rE_{i4},&
 W_{ij}(r)&\longmapsto I_5+t^rE_{ij}.
\end{aligned}
\label{eq:finite-matrix-images}
\end{equation}
Every displayed elementary matrix is invertible, with
$(I_5+aE_{ij})^{-1}=I_5-aE_{ij}$ for $i\ne j$.
Also $(I_5+aE_{ij})^p=I_5$ in characteristic $p$.

We check the relations.  Diagonal conjugation multiplies $E_{ab}$
by the ratio of the diagonal entries at $a$ and $b$, which gives
exactly \eqref{eq:group-weights}.  The identity
$E_{ab}E_{cd}=\delta_{bc}E_{ad}$ shows that all products in both
orders vanish for each required commuting pair: two units $E_{0i}$,
two units $E_{i4}$, the pair $E_{0i},E_{j4}$ for $i\neq j$, the pair
$E_{ij},E_{0h}$ for $h\neq i$, and the pair $E_{ij},E_{h4}$ for
$h\neq j$.  The required power relations hold by the same
square-zero calculation.  For $i\neq j$, the nonzero product in the
first noncommutation relation is $E_{0i}E_{ij}=E_{0j}$, and in the
second it is $E_{ij}E_{j4}=E_{i4}$.  The higher products vanish,
giving
\[
\begin{aligned}
 [I_5+t^rE_{0i},I_5+t^sE_{ij}]&=I_5+t^{r+s}E_{0j},\\
 [I_5+t^sE_{ij},I_5+t^rE_{j4}]&=I_5+t^{s+r}E_{i4}.
\end{aligned}
\]
Thus \eqref{eq:finite-matrix-images} defines a homomorphism from
$G$ into the finite group $\operatorname{GL}_5(S_m)$ of invertible
$5\times5$ matrices over $S_m$.  Let $F_m$
be its image.

The image of $z_N$ is $I_5+t^NE_{04}$.  It is central in $F_m$:
the unit $E_{04}$ has zero product in both orders with every
off-diagonal unit in \eqref{eq:finite-matrix-images}, and the
diagonal matrices have entry $1$ at both $0$ and $4$.  For
$c_0,\ldots,c_{m-1}\in\{0,\ldots,p-1\}$,
\[
 \prod_{q=0}^{m-1}(I_5+t^qE_{04})^{c_q}
   =I_5+\left(\sum_{q=0}^{m-1}c_qt^q\right)E_{04}.
\]
The classes $1,t,\ldots,t^{m-1}$ are an $\F_p$-basis of $S_m$,
because the defining polynomial is monic of degree $m$.
Thus this product is the identity only when all $c_q$ are zero.
The claimed independence and orders follow.
\end{proof}

\subsection{The algebra and its test modules}

\begin{proof}[Proof of \Cref{prop:selection}]
Take $R=kG$. A finite group presentation gives a finite algebra
presentation by adjoining inverses to its generators and imposing
the inverse equations and group relators. Extend $\alpha_G$ linearly
to $\alpha:R\to R$. For $q\in\Z$ define the central idempotents
\[
e_q=p^{-1}\sum_{a=0}^{p-1}z_q^a,\qquad e=e_0.
\]
Then $\alpha(e_q)=e_{q+1}$.
Use the same symbols for their images in $k[F_m]$, and put
\[
p_m=\left(\prod_{q=0}^{m-2}e_q\right)(1-e_{m-1}),
\]
with empty product $1$ when $m=1$.
The coefficient of the group basis element $z_{m-1}$ in $p_m$ is
$-p^{-m}\ne0$.
Indeed, the term $\prod_{q=0}^{m-2}e_q$ has no coefficient at
$z_{m-1}$, whereas the product with $e_{m-1}$ contributes $p^{-m}$
there. The independence in \Cref{lem:finite-central-quotients} makes
this contribution unique. Thus $k[F_m]p_m$ is nonzero.
For $0\le q<m-1$, the identities $e_qp_m=p_m$ and $e_{m-1}p_m=0$
show that $e_q$ acts as identity and $e_{m-1}$ as zero on this left
ideal and on each quotient, since these idempotents are central.
Choose a simple quotient $Y_m$. It is finite-dimensional and remains
simple under the surjection $R\to k[F_m]$.

For every $R$-module $Y$ and $j\ge0$, restriction of scalars gives
\begin{equation}\label{eq:selection-iterates}
H^{\circ j}(Y)\cong{}_{\alpha^j}(e_0e_1\cdots e_{j-1}Y).
\end{equation}
The product is $1$ for $j=0$. At stage $j$, selection by $e$ is the original action of
$\alpha^j(e)=e_j$, and the next pullback makes $r$ act as
$\alpha^{j+1}(r)$. This proves the formula by induction, and its
application to $Y_m$ proves \eqref{eq:extinction}.
\end{proof}

\section{Localization and a fixed lifted diagram}
\label{sec:localization}

Let $k$ be a field, let $R$ be a finitely presented unital $k$-algebra,
let $e\in R$ be a central idempotent, and let
$\alpha:R\to R$ be a unital $k$-algebra endomorphism.  Retain $H(Y)={}_{\alpha}(eY)$, now allowing $\alpha$ to be an
endomorphism. Centrality of $e$ still makes $eY$ stable under this action.

We first encode the presentation of $R$ in a finite-dimensional directed
algebra.  This is a three-vertex version of the homogenized quiver
construction in \cite[Section~1, Theorem~1.1 and Proposition~1.2]{NRS2004}.
The associated Verdier quotient carries the required $R$-action.

\subsection{A quadratic presentation and its directed encoding}

Choose a presentation
\[
  R=k\langle x_1,\ldots,x_d\rangle/(p_1,\ldots,p_t)
\]
in which each $p_j$ has degree at most two, allowing constant and linear
terms. Degree means the maximum word length among the nonzero monomials,
with the zero polynomial assigned degree $-\infty$.
Such a presentation exists: in any finite presentation, introduce
finitely many auxiliary generators for the successive partial products
in its finitely many monomials, together with their binary-product
definitions.  The resulting relations have degree at most two, and
eliminating the auxiliary generators recovers the original presentation.

Consider the quiver with vertices $0,1,2$ and arrow sets
\[
  \mathcal A_{01}=\{s,a_1,\ldots,a_d\},\qquad
  \mathcal A_{12}=\{s',a'_1,\ldots,a'_d\},\qquad
  \mathcal A=\mathcal A_{01}\sqcup\mathcal A_{12},
\]
where the subscripts specify the source and target. Impose the relations
\begin{equation}\label{eq:quadratic-encoding}
  a'_h s=s'a_h\quad(1\le h\le d),
  \qquad \widetilde p_j=0\quad(1\le j\le t),
\end{equation}
where homogenization replaces
\[
  x_hx_l\longmapsto a'_h a_l,\qquad
  x_h\longmapsto s'a_h,\qquad
  1\longmapsto s's.
\]
Let $B$ be the resulting path algebra with relations, and let
$\mathbf e_0,\mathbf e_1,\mathbf e_2$ be its vertex idempotents.
Write $\mathcal R$ for a fixed vector-space basis of the span of all
the relations in \eqref{eq:quadratic-encoding}, viewed in the free
space of paths of length two.  Thus each $\rho\in\mathcal R$ has a
unique expression
\[
  \rho=\sum_{\substack{a\in\mathcal A_{01}\\b\in\mathcal A_{12}}}
             \rho_{ba}\,ba.
\]
The basis $\mathcal R$ generates the same relation ideal.  There are
no paths of length three, so $B$ is finite-dimensional.

The global dimension $\gldim\Lambda$ is the supremum of the
projective dimensions of all left $\Lambda$-modules; the right global
dimension is $\gldim(\Lambda^{\mathrm{op}})$, where
$\Lambda^{\mathrm{op}}$ is the algebra with reversed multiplication.

\begin{lemma}\label{lem:directed-bound}
Let $\Lambda$ be a finite-dimensional $k$-algebra with complete
orthogonal idempotents $\varepsilon_0,\ldots,\varepsilon_{n-1}$ such that
\[
  1=\sum_{i=0}^{n-1}\varepsilon_i,\qquad
  \varepsilon_i\Lambda\varepsilon_i=k\varepsilon_i,\qquad
  \varepsilon_j\Lambda\varepsilon_i=0\quad(j<i).
\]
Then $\Lambda$ has global dimension at most $n-1$ on both sides.
In particular, the algebra $B$ above has global dimension at most two
on both sides.
\end{lemma}

\begin{proof}
Say that a left module $M$ is supported at indices at least $p$ if
$\varepsilon_jM=0$ for $j<p$.  There is a projective surjection
\[
  \bigoplus_{i=p}^{n-1}
      \Lambda\varepsilon_i\otimes_k\varepsilon_iM
      \longrightarrow M,
  \qquad c\otimes v\longmapsto cv.
\]
The source is projective even when $M$ is not finitely generated,
since each summand is a direct sum of copies of
$\Lambda\varepsilon_i$.  Its components below $p$ vanish, and its
component at $p$ maps isomorphically onto $\varepsilon_pM$:
only the summand $i=p$ contributes there.  Its kernel is therefore
supported at indices at least $p+1$.

Starting with $p=0$ and repeating this construction, the
$(n-1)$st kernel is supported only at the last vertex.  Every such
module is projective, since
$\Lambda\varepsilon_{n-1}=k\varepsilon_{n-1}$.
This gives the left dimension bound, including the case $n=1$.
The same proof for $\Lambda^{\mathrm{op}}$, with the reverse ordering
of the idempotents, gives the right bound.  The path directions and
the homogeneous quadratic relations of $B$ give precisely the
displayed hypotheses with $n=3$.
\end{proof}

For a finite-dimensional left $R$-module $Y$, define $M(Y)$ to be the
left $B$-module with $Y$ at every vertex, with $s,s'$ acting as the
identity, and with $a_h,a'_h$ acting as $x_h$.  The relations of $B$
hold because they evaluate to the defining relations of $R$.

\subsection{The quotient action and evaluation}

Let
\[
  K=\K^b(\proj\text{-}B),\qquad E_i=\mathbf e_iB,
\]
where $i=0,1,2$. Here $K$ is the homotopy category of bounded complexes of finitely
generated projective \emph{right} $B$-modules.  With our path convention,
an arrow $a:i\to j$ belongs to $\mathbf e_jB\mathbf e_i$.
Left multiplication by $a$ is consequently a right-module map
$E_i\to E_j$.  Composing these maps agrees with multiplying paths
in the stated order.

For a chain map $u:T\to T'$, its mapping cone has terms
$\Cone(u)^n=(T')^n\oplus T^{n+1}$ and differential
$(y,x)\mapsto(d_{T'}y+u(x),-d_Tx)$. Its isomorphism class
in $K$ depends only on the homotopy class of $u$.
A thick subcategory is a full triangulated subcategory closed under
direct summands. Let $\mathcal S$ be the thick subcategory of $K$ generated by the
cones of $s:E_0\to E_1$ and $s':E_1\to E_2$, and form
\[
  q:K\longrightarrow Q=K/\mathcal S.
\]
We use the same notation for objects and arrows after applying $q$.
The maps $s$ and $s'$ are isomorphisms in $Q$.  Set
\[
  \overline x_h=s^{-1}a_h\in\End_Q(E_0).
\]
The relation $a'_h s=s'a_h$ gives $a'_h=s'a_hs^{-1}$, whence
\[
  (s's)^{-1}a'_ha_l
     =s^{-1}a_hs^{-1}a_l
     =\overline x_h\overline x_l.
\]
Likewise $(s's)^{-1}s'a_h=\overline x_h$ and
$(s's)^{-1}s's=\id_{E_0}$.  The homogenized relations therefore give
a unital homomorphism
\begin{equation}\label{eq:quotient-action}
  \theta:R\longrightarrow\End_Q(E_0),
  \qquad x_h\longmapsto\overline x_h.
\end{equation}

For each finite-dimensional $Y$, tensoring defines an exact functor
\[
  \operatorname{ev}_Y:K\longrightarrow\D^b(k),
  \qquad T\longmapsto T\otimes_B M(Y).
\]
Here $\D^b(k)$ denotes the bounded derived category of finite-dimensional
$k$-vector spaces; termwise projectivity of $T$ computes the derived
tensor product.  Under the natural identifications
\[
  E_i\otimes_B M(Y)\cong\mathbf e_iM(Y)=Y,
\]
an arrow map evaluates to its action on $M(Y)$.  Thus
$\operatorname{ev}_Y$ takes $s,s'$ to identities and their cones to zero.
The kernel of this exact functor is thick, so it contains $\mathcal S$.
The universal property of the Verdier quotient
\cite[Tag 05RJ]{StacksProject} therefore gives a factorization through $Q$.
We retain the notation $\operatorname{ev}_Y$ for this factorization.
On $E_0$ the homomorphism \eqref{eq:quotient-action} evaluates to the
given $R$-action on $Y$.

\subsection{An odd double of a formal summand}

The idempotent $\theta(e)$ evaluates to the projection of $Y$ onto
$eY$, but it need not split inside $Q$. We therefore seek an object
of $Q$ whose evaluation is $eY$ plus a shifted copy. The following
lemma gives such an object for a copy shifted by three.

For an additive category $\mathcal T$, its idempotent completion
$\Kar(\mathcal T)$ has objects $(T,p)$, where $p^2=p\in\End(T)$,
and morphisms $u:(T,p)\to(T',p')$ satisfying $u=p'up$.
The identity of $(T,p)$ is $p$, and the embedding
$T\mapsto(T,\id_T)$ is fully faithful.
If $\mathcal T$ has a shift functor, it extends
by $(T,p)[1]=(T[1],p[1])$.

For the general triangulated structure on the idempotent completion,
see \cite[Theorem~1.5]{BalmerSchlichting2001}. The following argument
uses triangles in $\mathcal T$.

\begin{lemma}\label{lem:odd-double}
Let $\mathcal T$ be a triangulated category, let $L\in\mathcal T$,
and let $\epsilon\in\End_{\mathcal T}(L)$ be idempotent.
Put $U=(L,\epsilon)\in\Kar(\mathcal T)$.
Both $U\oplus U[1]$ and $U\oplus U[3]$ are isomorphic in
$\Kar(\mathcal T)$ to objects of $\mathcal T$.
\end{lemma}

\begin{proof}
First consider a triangle from $\mathcal T$,
\[
  A\xrightarrow{f}B\xrightarrow{g}T\xrightarrow{h}A[1],
\]
and suppose that formal decompositions
$A=I\oplus A_0$, $B=I\oplus B_0$ identify $f$ with the identity
on $I$ and zero elsewhere. Write $\operatorname{pr}_{B_0}$ for
the projection $B\to B_0$ and $\operatorname{in}_{A_0[1]}$ for the
inclusion $A_0[1]\to A[1]$. For any $Z\in\Kar(\mathcal T)$,
the long exact sequence $\Hom(Z,-)$ still holds for this triangle:
$Z$ is a retract of an object of $\mathcal T$, and its Hom sequence
is a direct summand of the exact sequence for that object.

The equalities $gf=0$ and $f[1]h=0$ give maps $g_0:B_0\to T$
and $h_0:T\to A_0[1]$ with
\[
  g=g_0\operatorname{pr}_{B_0},\qquad
  h=\operatorname{in}_{A_0[1]}h_0.
\]
Writing $(g_0)_*$ and $(h_0)_*$ for postcomposition by these maps,
exactness gives, for every such $Z$, a short exact sequence
\begin{equation}\label{eq:formal-cone-hom}
  0\longrightarrow\Hom(Z,B_0)
   \xrightarrow{(g_0)_*}\Hom(Z,T)
   \xrightarrow{(h_0)_*}\Hom(Z,A_0[1])
   \longrightarrow0.
\end{equation}
Indeed, a map into $B_0$ killed by $g_0$ factors through $f$ after
inclusion into $B$, and projection back to $B_0$ makes it zero.
At the other end, every map into $A_0[1]$ is killed by $f[1]$ after
inclusion into $A[1]$, so it lifts through $h$.
Taking $Z=A_0[1]$ supplies a section $\sigma$ of $h_0$.
It follows from \eqref{eq:formal-cone-hom} that
\[
  (g_0,\sigma):B_0\oplus A_0[1]\longrightarrow T
\]
induces isomorphisms on $\Hom(Z,-)$ for every $Z$. Explicitly, a map
$t:Z\to T$ is the sum of $\sigma h_0t$ and a unique map factoring
through $g_0$, by \eqref{eq:formal-cone-hom}. The same exact sequence
shows uniqueness of this decomposition. Taking $Z=T$ gives a right
inverse to $(g_0,\sigma)$; injectivity with $Z=B_0\oplus A_0[1]$
shows that this inverse is also a left inverse.

Apply this calculation to a triangle on $1-\epsilon:L\to L$.
Its identity summand is $(L,1-\epsilon)$, and its zero summand on
both sides is $U$.  Its cone is an object $S\in\mathcal T$ with
$S\cong U\oplus U[1]$ in the completion.  Under this identification,
consider the map
\[
  S[1]\cong U[1]\oplus U[2]
     \xrightarrow{\left(\begin{smallmatrix}0&0\\\id&0\end{smallmatrix}\right)}
  U\oplus U[1]\cong S.
\]
Fullness of $\mathcal T\hookrightarrow\Kar(\mathcal T)$ makes
this a morphism between the old objects $S[1]$ and $S$.
Its cone is again in $\mathcal T$.
The preceding calculation, with common identity summand $U[1]$,
source complement $U[2]$, and target complement $U$, identifies
that cone with $U\oplus U[3]$.
\end{proof}

Apply \Cref{lem:odd-double} with $\mathcal T=Q$, $L=E_0$, and
$\epsilon=\theta(e)$, and write
\[
  U=(E_0,\theta(e)).
\]
Choose $V_0\in Q$ and an isomorphism $V_0\cong U\oplus U[3]$
in $\Kar(Q)$ as in \Cref{lem:odd-double}.
Centrality of $e$ gives a unital action of $R$ on $U$ by
\[
  r\longmapsto\theta(e\alpha(r)e)\in\End_{\Kar(Q)}(U).
\]
Its unit is $\theta(e)=\id_U$, and its products agree with those
in $R$ because $e$ commutes with every $\alpha(r)$.
Use the same action on $U[3]$ and transport the resulting action
on $U\oplus U[3]$ to $V_0$.

A \emph{$B$-diagram} in an additive $k$-linear category consists of
objects at the three vertices and arrow maps satisfying the relations
of $B$. The action above gives such a diagram in $Q$: place $V_0$ at
each vertex, let $s,s'$ act as identities, and let $a_h,a'_h$ act
by the endomorphism associated to $x_h$.

The category $\D^b(k)$ splits idempotents: every bounded complex of
vector spaces is isomorphic to its cohomology with zero differential,
where splitting can be done degree by degree.  Thus evaluation extends
additively to $\Kar(Q)$.  It takes $U$ to $eY$, and the preceding
action of $r$ to the restriction of $\alpha(r)$ to $eY$.
Consequently the evaluated $B$-diagram has
\begin{equation}\label{eq:target-evaluation}
  eY\oplus(eY)[3]
\end{equation}
at every vertex, identities on $s,s'$, and the action of
$\alpha(x_h)$ on both summands for $a_h,a'_h$.

\subsection{Clearing denominators for the finite diagram}

To represent this quotient diagram in the homotopy category, use
the denominator class
\[
  \Sigma=\{u:T\to T'\text{ in }K:\Cone(u)\in\mathcal S\}.
\]
By \cite[Tag 05RG]{StacksProject}, $\Sigma$ is a multiplicative
system compatible with the triangulated structure, and $Q$ is its
Verdier localization. The right-fraction rules used below also follow
directly from triangles.
Identities belong to $\Sigma$, and composition preserves $\Sigma$
by the octahedral axiom.  If $u:Y'\to Y$ belongs to $\Sigma$ and
$f:X\to Y$ is any map, complete $u$ to a triangle
\[
  Y'\xrightarrow{u}Y\xrightarrow{c}T\longrightarrow Y'[1],
  \qquad T\in\mathcal S.
\]
Complete $cf$ to a triangle
$X'\xrightarrow{v}X\xrightarrow{cf}T\to X'[1]$.
Then $v\in\Sigma$, and exactness of $\Hom(X',-)$ gives
$g:X'\to Y'$ with $fv=ug$.  This is the required Ore square.

For the cancellation rule, suppose that $u:Y\to Y'$ belongs to
$\Sigma$ and $ug=0$ for $g:X\to Y$.
A triangle $T[-1]\to Y\xrightarrow{u}Y'\to T$ shows that $g$
factors through a map $b:X\to T[-1]$ with $T\in\mathcal S$.
Completing $b$ to a triangle gives a denominator $v:X'\to X$
with $bv=0$, hence $gv=0$.  Applying this to differences gives
the corresponding rule for equality of morphisms.

These rules construct the localization by roofs
\[
  X\xleftarrow{u}X'\xrightarrow{f}Y,
  \qquad u\in\Sigma,
\]
representing $q(f)q(u)^{-1}$; equality is tested after a common source
refinement. In particular, every morphism $q(X)\to q(Y)$ has this
form, and
\begin{equation}\label{eq:zero-fraction}
  q(g)=0\quad\Longrightarrow\quad
  gv=0\text{ for some }v:X'\to X\text{ in }\Sigma
\end{equation}
for $g:X\to Y$. For the last assertion, compare the roofs
$(\id_X,g)$ and $(\id_X,0)$ using \cite[Tag 04VJ]{StacksProject}.
Finite families can be treated with one denominator by collecting
their maps into a single map to a finite direct sum of targets
\cite[Tag 04VI]{StacksProject}.

\begin{lemma}\label{lem:diagram-lift}
Every $B$-diagram in $Q$ is isomorphic to the image under $q$ of a
$B$-diagram in $K$.
\end{lemma}

\begin{proof}
Write the given objects as $W_i$ and its arrow maps as $t_a$.
Choose $C_i\in K$ and isomorphisms $\psi_i:q(C_i)\to W_i$.
First collect the maps
\[
  (\psi_2^{-1}t_b\psi_1)_{b\in\mathcal A_{12}}:
     q(C_1)\longrightarrow
     \bigoplus_{b\in\mathcal A_{12}}q(C_2).
\]
Represent this tuple by one roof, with denominator
$u_1:D_1\to C_1$ and numerator components $f_b:D_1\to C_2$.
Set $D_2=C_2$ and
\[
  \phi_1=\psi_1q(u_1),\qquad \phi_2=\psi_2.
\]
Then $\phi_2q(f_b)=t_b\phi_1$ for every $b\in\mathcal A_{12}$.

Next collect the adjusted maps
\[
  (\phi_1^{-1}t_a\psi_0)_{a\in\mathcal A_{01}}:
     q(C_0)\longrightarrow
     \bigoplus_{a\in\mathcal A_{01}}q(D_1).
\]
A single roof gives a denominator $u_0:D_0\to C_0$ and maps
$f_a:D_0\to D_1$.  Put $\phi_0=\psi_0q(u_0)$.
All the arrow squares now commute:
\begin{equation}\label{eq:lift-arrow-squares}
  \phi_jq(f_a)=t_a\phi_i\qquad(a:i\to j).
\end{equation}

For $\rho\in\mathcal R$, form the morphism in $K$
\[
  g_\rho=\sum_{b,a}\rho_{ba}f_bf_a:D_0\longrightarrow D_2.
\]
The relations for $t_a$ and \eqref{eq:lift-arrow-squares} give
$\phi_2q(g_\rho)=0$, so $q(g_\rho)=0$.
Apply \eqref{eq:zero-fraction} to the tuple $(g_\rho)_\rho$.
There is one denominator $v:D'_0\to D_0$ such that
$g_\rho v=0$ in $K$ for every $\rho$.
Replace $D_0$ by $D'_0$, every map $f_a$ with source vertex $0$ by
$f_av$, and $\phi_0$ by $\phi_0q(v)$.
The arrow squares still commute, and each relation becomes
\[
  \sum_{b,a}\rho_{ba}f_b(f_av)=g_\rho v=0
  \quad\text{in }K.
\]
There are no arrows into vertex zero and no longer paths requiring
further compatibility.  Thus the new data form a $B$-diagram in $K$,
and the three $\phi_i$ give an isomorphism from its image to the
entire prescribed diagram in $Q$.
\end{proof}

\subsection{The chain data for tensor realization}

The lifted relations are null-homotopic. Their homotopies give the
chain data for tensor realization.

\begin{proposition}[Fixed lifted diagram]\label{prop:lifted-diagram}
For the fixed data $R,e,\alpha$ and the algebra $B$ above, there are
bounded complexes $D_0,D_1,D_2$ of finitely generated projective
right $B$-modules, chain maps
\[
  f_a:D_i\longrightarrow D_j\qquad(a:i\to j),
\]
and right-$B$-linear maps $h_\rho:D_0\to D_2$ of degree $-1$
for $\rho\in\mathcal R$, such that
\begin{equation}\label{eq:relation-homotopies}
  d_{D_2}h_\rho+h_\rho d_{D_0}
      =\sum_{b,a}\rho_{ba}f_bf_a.
\end{equation}
For every finite-dimensional left $R$-module $Y$, the diagram
whose objects and arrows in $\D^b(k)$ are
\begin{equation}\label{eq:evaluated-diagram}
  D_i\otimes_B M(Y),\qquad f_a\otimes_B\id_{M(Y)}
\end{equation}
is isomorphic to the diagram having $eY\oplus(eY)[3]$ at all three
vertices, identities on $s,s'$, and the action of $\alpha(x_h)$
on both summands for $a_h,a'_h$.
All the complexes, maps, and homotopies in this statement are fixed
before $Y$ is chosen.
\end{proposition}

\begin{proof}
Apply \Cref{lem:diagram-lift} to the diagram on $V_0$. Choose bounded
right-projective complex representatives $D_i$ and chain-map
representatives $f_a$.  Each relation is zero in $K$, hence its
representative is null-homotopic.  Choose one right-linear
homotopy $h_\rho$ for each of the finitely many relations.
This gives \eqref{eq:relation-homotopies}.

Apply $\operatorname{ev}_Y$ to the arrow squares \eqref{eq:lift-arrow-squares}
and the identification $V_0\cong U\oplus U[3]$. Equation
\eqref{eq:target-evaluation} gives the required diagram isomorphism.

\end{proof}

\section{Realization by a bimodule complex}
\label{sec:realization}

We retain the algebra \(B\), its relation basis \(\mathcal R\), and the
module \(M(Y)\) from the preceding section.

\begin{theorem}\label{thm:realization}
Let \(k\) be a field, let \(R\) be a finitely presented \(k\)-algebra,
let \(e\in R\) be a central idempotent, and let
\(\alpha:R\to R\) be a unital $k$-algebra endomorphism.
Use the finite-dimensional algebra \(B\) and the modules \(M(Y)\)
constructed in the preceding section, and put
\(H(Y)={}_{\alpha}(eY)\).
There is a single bounded complex \(P\) of finite-dimensional
\(B\)-bimodules, termwise projective over $B^e=B\otimes_k B^{\mathrm{op}}$, such that for
every finite-dimensional left \(R\)-module \(Y\) one has
\begin{equation}\label{eq:realization}
 P\otL_B M(Y)
 \simeq M(H(Y))\oplus M(H(Y))[3]
 \qquad\text{in }\D^b(B\text{-}\modu).
\end{equation}
\end{theorem}

To prove the theorem, we first replace the diagram of
\Cref{prop:lifted-diagram} by a complex of \(B\)-bimodules.
The three-column construction below incorporates the relation homotopies;
compare the general homotopy-action framework in \cite{Keller2000}.

\begin{lemma}\label{lem:rectification}
Let \(D_0,D_1,D_2\) be bounded complexes of finitely generated projective
right \(B\)-modules.  Suppose that each arrow \(a:i\to j\) of \(B\)
is assigned a chain map \(f_a:D_i\to D_j\), and that each relation
\[
 \rho=\sum_{b,a}\rho_{ba}\,ba\in\mathcal R
\]
is assigned a right \(B\)-linear map \(h_\rho:D_0\to D_2\) of degree
\(-1\) satisfying
\[
 d_{D_2}h_\rho+h_\rho d_{D_0}
 = \sum_{b,a}\rho_{ba}f_bf_a .
\]
There is a bounded complex \(P\) of finite-dimensional \(B\)-bimodules,
termwise projective as a right \(B\)-module, together with homotopy
equivalences of right-module complexes
\[
 \iota_i:D_i\longrightarrow \mathbf e_iP \qquad (i=0,1,2),
\]
such that, for every arrow \(a:i\to j\), the two maps
\(a\iota_i\) and \(\iota_jf_a\) from \(D_i\) to \(\mathbf e_jP\)
are chain homotopic.
\end{lemma}

\begin{proof}
All tensor products in this construction are over \(k\). Form the
bimodule complexes
\begin{align*}
 L'_0&=\bigoplus_{i=0}^2 B\mathbf e_i\otimes_k D_i,\\
 L'_1&=\bigoplus_{a:i\to j}B\mathbf e_j\otimes_k D_i,\\
 L'_2&=\bigoplus_{\rho\in\mathcal R}B\mathbf e_2\otimes_k D_0.
\end{align*}
The left \(B\)-action is on the first tensor factor and the right
\(B\)-action is on the second.  Let
\(d_0,d_1,d_2\) denote the respective internal differentials of these
three complexes.

Define $p:L'_1\to L'_0$ on the summand indexed by \(a:i\to j\) by
\[
 p(c\otimes v)
  =(ca\otimes v)_i-(c\otimes f_a(v))_j .
\]
Here subscripts indicate summands of \(L'_0\). Define
$q:L'_2\to L'_1$ on the summand indexed by \(\rho\) by
\[
 q(c\otimes v)
  =\sum_{b,a}\rho_{ba}
       \bigl((cb\otimes v)_a+(c\otimes f_a(v))_b\bigr).
\]
These are bimodule chain maps of degree zero.  Define also the
bimodule map \(h:L'_2\to L'_0\) of degree \(-1\) by
\[
 h(c\otimes v)=(c\otimes h_\rho(v))_2
\]
on the relation summand \(\rho\).  Expanding the composite \(pq\)
gives
\begin{align*}
 pq(c\otimes v)
 &=\sum_{b,a}\rho_{ba}
   \bigl(
     (cba\otimes v)_0-(cb\otimes f_a(v))_1\\
 &\hspace{40mm}
     +(cb\otimes f_a(v))_1-(c\otimes f_bf_a(v))_2
   \bigr)\\
 &=-\sum_{b,a}\rho_{ba}(c\otimes f_bf_a(v))_2.
\end{align*}
The last equality uses \(\rho=0\) in \(B\).

Set
\[
 P=L'_0\oplus L'_1[1]\oplus L'_2[2]
\]
as a graded bimodule, and equip it with the differential
\[
 d_P=
 \begin{pmatrix}
   d_0 & p    & h\\
   0   & -d_1 & q\\
   0   & 0    & d_2
 \end{pmatrix}.
\]
For example, \(P^n=(L'_0)^n\oplus(L'_1)^{n+1}
\oplus(L'_2)^{n+2}\), so \(p,q,h\) each have total degree one
in this matrix.  The diagonal signs agree with our shift convention.
Since \(p\) and \(q\) commute with the internal differentials,
the \((0,1)\) and \((1,2)\) entries of \(d_P^2\) vanish.
Its remaining off-diagonal entry is
\[
 d_0h+pq+hd_2=0,
\]
by the displayed calculation and the identity for \(h_\rho\).
The diagonal entries vanish as well.  Thus \(P\) is a complex of
bimodules.

Every term \(B\mathbf e_j\otimes_k D_i^n\), viewed as a right
\(B\)-module, is a finite direct sum of copies of \(D_i^n\).
Each $D_i^n$ is a direct summand of a finite free right $B$-module.
Tensoring its splitting with $B\mathbf e_j$ exhibits
$B\mathbf e_j\otimes_kD_i^n$ as a direct summand of finitely many
copies of $B\mathbf e_j\otimes_kB$. As left $B^e$-modules,
\[
 B^e(\mathbf e_j\otimes1)
 \xrightarrow{\ \sim\ }B\mathbf e_j\otimes_kB,
 \qquad b\mathbf e_j\otimes c^{\mathrm{op}}
 \longmapsto b\mathbf e_j\otimes c.
\]
Here $(a\otimes d^{\mathrm{op}})(v\otimes c)=av\otimes cd$,
so the displayed map respects the action. The left-hand side is an
idempotent summand of $B^e$, proving the required $B^e$-projectivity. All terms of $P$ are
finite dimensional and projective on both sides.  There are finitely many summands, and the complexes
\(D_i\) are bounded; hence \(P\) is bounded.

Define
\[
 \iota_i(v)=(\mathbf e_i\otimes v)_i
       \in \mathbf e_iL'_0\subseteq \mathbf e_iP .
\]
This is a chain map and is split injective as a map of graded right
modules.  We prove that its quotient is acyclic.

Call \(j\) the \emph{source index} of a summand whose second tensor
factor is \(D_j\), in any of the three columns.  Give
\(\mathbf e_iP\) the decreasing filtration
\[
 F^r(\mathbf e_iP)
  =\text{the sum of its summands with source index at least }r
  \qquad (0\leq r\leq 3).
\]
These are subcomplexes.  Indeed, the internal differentials and the
path-multiplication parts of \(p,q\) preserve the source index;
the \(f_a\) terms increase it; and \(h\) increases it from \(0\)
to \(2\).  In particular \(F^0=\mathbf e_iP\) and \(F^3=0\).
The quotient by \(\iota_i(D_i)\) inherits this finite filtration.

On the associated graded, only the internal differentials and the
path-multiplication parts remain.  The source-\(i\) part consists
solely of \(\iota_i(D_i)\), since the only path from \(i\) to itself
is the trivial path.  It disappears in the quotient.  Source indices
larger than \(i\) give zero. For a finite set $\Omega$, let
$k^{\Omega}$ denote the vector space with basis indexed by $\Omega$.
If \(j<i\) are adjacent vertices, the
source-\(j\) part is the signed total complex obtained by tensoring
\(D_j\) with the exact sequence
\[
 0\longrightarrow
 k^{\{a:j\to i\}}
 \xrightarrow{\ a\mapsto a\ }
 \mathbf e_iB\mathbf e_j
 \longrightarrow 0.
\]
There are no relations of length one, so the displayed map is an
isomorphism.

For the remaining pair \((j,i)=(0,2)\), that part is the signed total
complex of \(D_0\) tensored with
\[
 0\longrightarrow k^{\mathcal R}
 \longrightarrow k^{\{(b,a):a:0\to1,\ b:1\to2\}}
 \longrightarrow \mathbf e_2B\mathbf e_0
 \longrightarrow 0.
\]
The three nonzero horizontal terms have degrees \(-2,-1,0\).
The first map sends the basis vector indexed by \(\rho\) to
\(\sum_{b,a}\rho_{ba}(b,a)\); the second sends \((b,a)\) to \(ba\).
The first map is injective because \(\mathcal R\) is a basis of the
relation space.  Its image is the kernel of the second map: all
defining relations have length two, and no path of positive length
can enter vertex \(0\) or leave vertex \(2\).  Thus the ideal they
generate, in this path space, is exactly their linear span.  The
second map is surjective because length-two paths span
\(\mathbf e_2B\mathbf e_0\).

These exact complexes of vector spaces are split, and tensoring
with a bounded complex \(D_j\) gives an acyclic total complex.
All associated graded pieces of the quotient are therefore acyclic.
Induction through the finite filtration proves that the quotient
itself is acyclic.  Thus \(\iota_i\) is a quasi-isomorphism.
Its cone is a bounded acyclic complex of projective right modules,
and is contractible: starting in the largest nonzero degree, split
the surjection onto its projective last term and continue downward.
It follows that \(\iota_i\) is a homotopy equivalence.

Finally, for \(a:i\to j\) define a right-linear map of degree \(-1\)
\[
 H_a:D_i\longrightarrow\mathbf e_jP,\qquad
 H_a(v)=(\mathbf e_j\otimes v)_a
          \in\mathbf e_jL'_1[1].
\]
The internal differential on this summand is \(-d_{D_i}\);
its contribution cancels that of \(H_ad_{D_i}\).  The remaining
component is \(p\), giving
\[
 d_PH_a+H_ad_{D_i}
   =a\iota_i-\iota_jf_a .
\]
This is the required arrow homotopy.
\end{proof}

\begin{proof}[Proof of \Cref{thm:realization}]
Choose the fixed complexes \(D_i\), arrow maps \(f_a\), and relation
homotopies \(h_\rho\) of \Cref{prop:lifted-diagram}, and apply
\Cref{lem:rectification}.  Since \(P\) and the \(D_i\) are bounded
complexes of projective right modules, their ordinary tensor
products with \(M(Y)\) compute the derived tensor products.
If chain maps $u,v$ satisfy $u-v=dh+hd$, tensoring this identity
with $M(Y)$ gives the same identity with homotopy $h\otimes\id$.
Applying this to the homotopy inverses of $\iota_i$ and to the arrow
homotopies preserves both comparisons after tensoring.  Thus the vertex diagram of
\[
 T_Y:=P\otL_B M(Y)
\]
in the derived category of \(k\)-vector spaces is isomorphic to the
evaluated diagram of \Cref{prop:lifted-diagram}.

At each vertex this diagram is \(eY\oplus eY[3]\).  The arrows
\(s,s'\) act by identity maps, and \(a_h,a'_h\) act by
\(\alpha(x_h)\) on each summand.  Applying cohomology, with its
induced arrow maps, therefore gives isomorphisms of left
\(B\)-modules
\[
 H^n(T_Y)\cong
 \begin{cases}
  M(H(Y)),& n=0,-3,\\
  0,& n\ne0,-3.
 \end{cases}
\]


Using these cohomology identifications, the canonical truncation
triangle is
\[
 M(H(Y))[3]\longrightarrow T_Y\longrightarrow M(H(Y))
 \longrightarrow M(H(Y))[4].
\]
The connecting morphism belongs to
\[
 \Hom_{\D^b(B\text{-}\modu)}(M(H(Y)),M(H(Y))[4])
   =\Ext_B^4(M(H(Y)),M(H(Y)))=0,
\]
where the vanishing follows from the bound
\(\gldim B\leq2\) in \Cref{lem:directed-bound}.
Exactness of $\Hom(M(H(Y)),-)$ gives a section of the map
$T_Y\to M(H(Y))$, since the connecting morphism is zero. Together
with the first arrow, this section gives
$M(H(Y))[3]\oplus M(H(Y))\to T_Y$. It induces isomorphisms in
degrees $-3$ and $0$, and all other cohomology groups vanish.
It is therefore an isomorphism in the derived category, proving
\eqref{eq:realization}.

\end{proof}

\begin{corollary}
Put \(\Phi=P\otL_B-\).  For every integer \(j\geq0\) and every
finite-dimensional left \(R\)-module \(Y\), there is an isomorphism
\[
 \Phi^j(M(Y))
 \simeq
 \bigoplus_{q=0}^j
   \bigl(M(H^{\circ j}(Y))[3q]\bigr)^{\oplus\binom{j}{q}} .
\]
\end{corollary}

\begin{proof}
The case \(j=0\) is immediate.  For the induction step, apply
\(\Phi\) to an isomorphism for \(j\); derived tensor preserves
isomorphisms, finite direct sums, and shifts.  Apply
\eqref{eq:realization} with \(H^{\circ j}(Y)\) in place of \(Y\) to each
summand. The summand indexed by $q$ contributes shifts $3q$ and
$3(q+1)$, each with multiplicity $\binom jq$. Thus the multiplicity
at shift $3q$ in the next iterate is
\[
 \binom jq+\binom j{q-1}=\binom{j+1}q,
\]
where $\binom jq=0$ for $q<0$ or $q>j$.
These nonnegative integers count copies in a direct sum.
Each $H^{\circ j}(Y)$ is finite-dimensional.

\end{proof}

\section{Simulation by an ordinary bimodule}
\label{sec:ordinary-simulation}

We encode the differential of $P$ in an ordinary bimodule. Two
derived tensor steps will recover $P\otL_B-$ up to a fixed shift.

Choose integers \(a_*\le 0\le b_*\) such that \(P^n=0\) outside
\([a_*,b_*]\), enlarging the interval if necessary, and put
\(l=b_*-a_*\).
Let \(C\) be the algebra of the quiver
\[
0\xrightarrow{c_1}1\xrightarrow{c_2}\cdots
\xrightarrow{c_l}l
\]
modulo all paths of length two.
Write \(\varepsilon_i\) for its vertex idempotents, and let \(W\)
be the simple \emph{right} \(C\)-module at vertex \(l\).
There is a projective resolution of right \(C\)-modules
\begin{equation}
\label{eq:right-chain-resolution}
0\longrightarrow \varepsilon_0C
\xrightarrow{c_1\cdot}\varepsilon_1C
\xrightarrow{c_2\cdot}\cdots
\xrightarrow{c_l\cdot}\varepsilon_lC
\longrightarrow W\longrightarrow0.
\end{equation}
Indeed, for \(i>0\), the right projective \(\varepsilon_iC\)
has basis \(\varepsilon_i,c_i\), and its radical is the right
simple at \(i-1\).  Left multiplication by \(c_i\) maps
\(\varepsilon_{i-1}\) to \(c_i\) and, when \(i>1\), maps
\(c_{i-1}\) to zero.  This proves exactness, including injectivity
at the first term.  For \(l=0\), the resolution consists of
the isomorphism \(\varepsilon_0C\to W\).

Define a \((B\otimes_k C,B)\)-bimodule \(O\) by
\[
\varepsilon_iO=P^{a_*+i},
\qquad
c_i\big|_{\varepsilon_{i-1}O}=d_P^{a_*+i-1}.
\]
Here the two \(B\)-actions are those on the terms of \(P\).
The relation \(d_P^2=0\) makes this a representation of \(C\),
and the bimodule linearity of \(d_P\) makes the \(C\)-action
commute with both \(B\)-actions.  Thus \(O\) is an ordinary
finite-dimensional bimodule.  Set
\[
D=B\times(B\otimes_k C),\qquad
T=B\otimes_k W,\qquad X=O\oplus T.
\]
Regard \(O\) as having its left action from the second factor
of \(D\) and its right action from the first.  Regard \(T\)
as having its left action from the first factor and its
right action from the second, with the regular \(B\)-actions
on its \(B\)-factor.  These prescriptions define an ordinary
finite-dimensional \(D\)-bimodule \(X\).

The two summands of \(X\) move complexes between the factors of \(D\),
as shown in \Cref{fig:simulation}. Thus two tensor steps return to the
first factor; the next proposition identifies their composite.

\begin{figure}[htbp]
\centering
\begin{tikzcd}[column sep=huge]
 \D^b(B\text{-}\modu)
 \arrow[r,bend left=22,"O\otL_B-"] &
 \D^b((B\otimes_k C)\text{-}\modu)
 \arrow[l,bend left=22,"T\otL_{B\otimes_k C}-"]
\end{tikzcd}
\caption{The two factors of \(D\) and the tensor functors supplied by
the off-diagonal summands of \(X\). Starting on the left, the composite
is \(P[b_*]\otL_B-\) by \Cref{prop:ordinary-simulation}.}
\label{fig:simulation}
\end{figure}

\begin{proposition}
\label{prop:ordinary-simulation}
The algebra \(D\) has finite global dimension on both sides.
Let \(F=X\otL_D-\).  For every complex \(Z\) of left
\(B\)-modules, regarded as a complex on the first factor of
\(D\), there is an isomorphism
\begin{equation}
\label{eq:compression}
F^2(Z)\simeq P[b_*]\otL_B Z,
\qquad
T\otL_{B\otimes_k C}O\simeq P[b_*].
\end{equation}
The second isomorphism is an isomorphism of complexes of
\(B\)-bimodules, and the first is supported on the first
factor of \(D\).
\end{proposition}

\begin{proof}
The vertices of \(B\otimes_k C\) are pairs \((i,j)\), where
\(0\le i\le2\) and \(0\le j\le l\).  Its diagonal vertex
corners are \(k\).  A nonzero component between distinct
vertices moves in the product of the two directed orders.
Assign height $i+j$ to vertex $(i,j)$. Every nonzero
off-diagonal component increases height; the algebra on each height
layer is a product of copies of $k$. If a left module $M$ is supported
in heights at least $r$, where $0\le r\le l+2$, multiplication gives
the projective surjection
\[
 \bigoplus_{i+j\ge r}(B\otimes_k C)
 (\mathbf e_i\otimes\varepsilon_j)\otimes_k
 (\mathbf e_i\otimes\varepsilon_j)M\longrightarrow M.
\]
At each vertex of height $r$, only its own summand contributes, and
the map is an isomorphism there. The kernel is therefore supported
in heights at least $r+1$. After $l+2$ steps it is supported at the
last height, where it is projective. Reversing the order proves the
right-module bound. Thus
\[
\gldim(B\otimes_k C),\quad
\gldim\bigl((B\otimes_k C)^{\mathrm{op}}\bigr)\le l+2.
\]
The same bound holds for $D$ and $D^{\mathrm{op}}$, since
$\gldim B,\gldim B^{\mathrm{op}}\le2$.

Tensor \eqref{eq:right-chain-resolution} over \(k\) with
the regular \(B\)-bimodule, placing its last projective term
in degree zero.  This gives a bounded resolution \(Q_T\to T\)
by right \(B\otimes_k C\)-projectives, preserving the left
\(B\)-action.  Also, \(O\) is projective as a right
\(B\)-module, being the finite direct sum of the terms of \(P\).
Thus \(Q_X=O\oplus Q_T\), with the prescribed block actions,
is a bounded right-\(D\)-projective resolution of \(X\) retaining
its left \(D\)-action.  Ordinary tensor with \(Q_X\) computes
\(F\) on every complex.  On the first factor, the block supports
therefore give
\[
F^2(Z)\simeq
Q_T\otimes_{B\otimes_k C}(O\otimes_B Z)
\cong
\bigl(Q_T\otimes_{B\otimes_k C}O\bigr)\otimes_B Z.
\]
In degree \(n\), for \(-l\le n\le0\), tensoring \(Q_T\)
with \(O\) gives
\[
\bigl(B\otimes_k\varepsilon_{l+n}C\bigr)
\otimes_{B\otimes_k C}O
\cong\varepsilon_{l+n}O
=P^{a_*+l+n}=P^{b_*+n}.
\]
The induced differential is \(d_P\).
In comparison, the standard differential of \(P[b_*]\) is
\((-1)^{b_*}d_P\).  Multiplication by
\((-1)^{b_*n}\) in degree \(n\) defines a chain isomorphism
to \(P[b_*]\); it respects both \(B\)-actions.
Since this complex is termwise right-\(B\)-projective, its
ordinary tensor with any \(Z\) computes the derived tensor.
This proves the second isomorphism in
\eqref{eq:compression}, and hence the first.
\end{proof}

\section{A square-zero extension and projective dimension}
\label{sec:square-zero}

For a $D$-bimodule $X$, the square-zero extension $D\ltimes X$
is the vector space $D\oplus X$ with multiplication
$(d,x)(d',x')=(dd',dx'+xd')$. A derived bar decomposition
relates tensor extinction to projective dimension.

\subsection{The derived bar decomposition}

The splitting below also follows from the graded-component decomposition
of Minamoto and Yamaura
\cite[Lemma~4.13(4) and the proof of Theorem~4.17]{MinamotoYamaura},
after forgetting the grading and passing to opposite algebras.
The direct proof identifies each derived tensor factor and its shift.

\begin{proposition}
\label{prop:bar-splitting}
Let \(D\) be a finite-dimensional \(k\)-algebra, let \(X\) be
a finite-dimensional \(D\)-bimodule, and let
\(A=D\ltimes X\).
Let \(N\) be a left \(D\)-module, inflated to \(A\) through
the projection \(A\to D\), and put \(F=X\otL_D-\).
In the unbounded derived category $\D(D\text{-}\Mod)$ of all left
$D$-modules there is an isomorphism
\begin{equation}
\label{eq:bar-splitting}
D\otL_A N
\simeq
\bigoplus_{r\ge0}(F^rN)[r].
\end{equation}
Here \(F^0N=N\), and every \(F^r\) denotes an iteration
of the derived tensor functor.
\end{proposition}

\begin{proof}
Use the ordinary unnormalized bar resolution of \(N\) over
the field \(k\).  It is a free left \(A\)-resolution with
degree \(-n\) term \(A\otimes_k A^{\otimes_k n}\otimes_k N\).
Its augmentation is a quasi-isomorphism: the augmented
bar complex has the usual \(k\)-linear contraction obtained
by inserting the unit.  After tensoring with the right
\(A\)-module \(D\), it yields a model \(\mathcal B\) for
\(D\otL_AN\), with
\[
\mathcal B^{-n}=D\otimes_k A^{\otimes_k n}\otimes_k N
\qquad(n\ge0).
\]
Write $\pi:A\to D$ for the projection. For $n\ge1$, the
differential on $d\otimes a_1\otimes\cdots\otimes a_n\otimes v$ is
\begin{align*}
 &d\pi(a_1)\otimes a_2\otimes\cdots\otimes a_n\otimes v\\
 &\quad+\sum_{i=1}^{n-1}(-1)^i
 d\otimes a_1\otimes\cdots\otimes(a_i a_{i+1})
       \otimes\cdots\otimes a_n\otimes v\\
 &\quad+(-1)^n d\otimes a_1\otimes\cdots\otimes a_{n-1}
       \otimes a_nv .
\end{align*}
For $n=1$ the middle sum is empty.

Expand each internal copy of \(A\) as \(D\oplus X\).
Every nonzero face preserves the number of entries from
\(X\).  Indeed, a product involving exactly one \(X\)-entry
remains in \(X\), a product of two \(X\)-entries is zero,
and an endpoint action by \(X\) is zero on both \(D\) and
\(N\). Thus there is a decomposition of complexes
\[
\mathcal B=\bigoplus_{r\ge0}\mathcal B_r,
\]
where \(\mathcal B_r\) consists of words with precisely
\(r\) entries from \(X\).

For fixed \(r\), write \(i_0,\ldots,i_r\) for the lengths
of the blocks of \(D\)-entries before, between, and after
the \(X\)-entries.  The corresponding tensor is
\begin{equation}
\label{eq:bar-block}
D\otimes_kD^{\otimes_k i_0}\otimes_k X
\otimes_kD^{\otimes_k i_1}\otimes_k\cdots
\otimes_k X\otimes_kD^{\otimes_k i_r}\otimes_k N,
\end{equation}
in degree \(-r-\sum_{j=0}^r i_j\).
We identify this with the multiple two-sided bar complex
for
\begin{equation}
\label{eq:derived-multibar}
D\otL_D
\underbrace{X\otL_D\cdots\otL_D X}_{r\ \text{factors}}
\otL_D N,
\end{equation}
shifted by \(r\).

First consider the signs.  Within block \(j\), a bar face
has alternating-sign offset
\[
j+\sum_{h<j}i_h.
\]
In the ordinary multibar total differential, ordered from
left to right, the corresponding offset is
\(\sum_{h<j}i_h\); shifting the total complex by \(r\)
adds the factor \((-1)^r\).
Rescale the component \eqref{eq:bar-block} by
\[
(-1)^{\sum_{j=0}^r(r-j)i_j}.
\]
When a face reduces \(i_j\) by one, this factor changes by
\((-1)^{r-j}\), precisely the discrepancy between the
two differentials.  This gives the required chain
isomorphism.  When \(r\ge1\) and \(i_j=0\), there is no
nonzero face in that block: the prospective face multiplies
two \(X\)-entries or applies an endpoint \(X\)-action.
For \(r=0\), this is the ordinary two-sided bar complex
for \(D\otL_DN\).  At \(i_0=0\) its differential is absent;
the map \(D\otimes_kN\to N\) is the separate augmentation.

To identify the multibar with \eqref{eq:derived-multibar},
we resolve each successive left $D$-complex by free modules.
For a bounded-above complex \(Z\) of left \(D\)-modules,
the total left-module bar resolution has terms obtained
from
\[
D\otimes_k D^{\otimes_k i}\otimes_k Z^q
\]
in degree \(q-i\).  Its terms are free left \(D\)-modules
because \(k\) is a field, and its augmentation to \(Z\)
is a quasi-isomorphism.  For each total degree there are
only finitely many possible \(i\), since \(Z\) is bounded
above.  Thus it is a bounded-above complex of free
modules and can be used to compute a derived tensor.

Starting at the right, resolve \(N\) in this way and
tensor with \(X\).  Resolve the resulting left
\(D\)-complex in the same way and tensor with the next
\(X\), continuing for the fixed finite number \(r\)
of factors.  A final bar resolution with the leading
factor \(D\) gives exactly the multibar
\eqref{eq:bar-block}, before its shift.
All exterior \(D\)-actions are preserved.
This establishes its identification with
\eqref{eq:derived-multibar}, which is isomorphic to
\(F^rN\).

Finally, in degree \(-n\) of \(\mathcal B\) only
\(0\le r\le n\) occur, and for each such \(r\) there
are finitely many tuples satisfying
\(r+\sum_j i_j=n\).
All totalizations above therefore use direct sums that
are finite in each degree.
Direct sums are exact for modules, so the componentwise
quasi-isomorphisms give \eqref{eq:bar-splitting}.
\end{proof}

\subsection{Detecting finite projective dimension and lower bounds}

\begin{corollary}
\label{cor:pd-detection}
In the setting of \Cref{prop:bar-splitting}, suppose that
\(D\) has finite global dimension on both sides and that
\(N\) is finite-dimensional.
\begin{enumerate}
\item If \(F^tN\simeq0\) for some \(t\ge1\), then
      \(\pd_A N<\infty\).
\item If \(F^rN\not\simeq0\), then \(\pd_A N\ge r\),
      with infinite projective dimension allowed.
\end{enumerate}
\end{corollary}

\begin{proof}
Every fixed iterate \(F^rN\) has bounded,
finite-dimensional cohomology.
For example, if \(g\) bounds the right global dimension
of \(D\), then for every finite-dimensional left
\(D\)-module \(V\) the groups
\(\Tor_i^D(X,V)\) are finite-dimensional and vanish for
\(i>g\).  Applying this observation to successive
truncation triangles shows inductively that
\[
H^q(F^rN)=0\quad\text{unless}\quad -rg\le q\le0.
\]

For the first assertion, vanishing at \(t\) implies
vanishing at every \(r\ge t\).
By \eqref{eq:bar-splitting},
\(D\otL_AN\) then has bounded cohomology.
Since \(X\) is a nilpotent ideal, it is contained in
\(\rad A\), and
\[
\rad A=(\rad D)\oplus X,
\qquad A/\rad A\cong D/\rad D=:S.
\]
Derived associativity gives
\[
S\otL_A N
\simeq S\otL_D(D\otL_A N).
\]
The right \(D\)-module \(S\) has a finite projective
resolution, so this complex also has bounded cohomology.

Choose a minimal projective resolution of \(N\) over
the finite-dimensional algebra \(A\), using successive
projective covers.  Its terms are finitely generated,
and every differential has image in the radical of
its target.  Tensoring with \(S=A/\rad A\) therefore
annihilates every differential.  If \(Q\) is a nonzero
finitely generated projective \(A\)-module, then
\[
S\otimes_A Q\cong Q/(\rad A)Q\ne0
\]
by Nakayama's Lemma.  Consequently bounded cohomology
of \(S\otL_AN\) forces the minimal resolution to have
only finitely many nonzero terms.  This proves
\(\pd_A N<\infty\).

For the second assertion, choose \(q\le0\) with
\(H^q(F^rN)\ne0\).  Its direct summand in
\eqref{eq:bar-splitting} gives
\[
H^{q-r}(D\otL_A N)\ne0,\qquad q-r\le-r.
\]
A projective \(A\)-resolution of \(N\) of length
strictly less than \(r\) would make the derived
tensor zero in this degree.  Hence
\(\pd_A N\ge r\), as asserted.
\end{proof}

With $D$ of finite global dimension, for a bounded complex $Z$
of left $D$-modules with finite-dimensional cohomology and $Z\not\simeq0$,
$\pd_D Z$ is the least integer $d$ such that $Z$ is represented by
a bounded projective complex zero below degree $-d$; set
$\pd_D0=-\infty$. In particular $\pd_D(Z[s])=\pd_D Z+s$.
\begin{proposition}\label{prop:exact-pd}
Under the hypotheses of \Cref{cor:pd-detection}, if $F^tN=0$ for
some $t\ge1$, then
\[
\pd_A N=\max_{0\le r<t}\{\pd_D(F^rN)+r\}.
\]
\end{proposition}
\begin{proof}
If $N=0$, then $F^rN=0$ for every $r\ge0$, and both sides are
$-\infty$ by our convention. Assume henceforth that $N\ne0$.
Put $S=D/\rad D=A/\rad A$. A minimal bounded projective complex
representing $Z$ has radical differentials, so
\[
\pd_D Z=\sup\{n\in\Z:H^{-n}(S\otL_D Z)\ne0\}.
\]
Such a complex exists by finite global dimension, followed by
splitting off contractible projective summands. Reduction modulo
the radical kills the remaining differentials and detects every
nonzero projective term by Nakayama's lemma.
By \Cref{cor:pd-detection}, $N$ has finite projective dimension.
Applying $S\otL_D-$ to \eqref{eq:bar-splitting} gives
\[
 S\otL_A N\simeq\bigoplus_{r=0}^{t-1}
 (S\otL_D F^rN)[r].
\]
For each surviving $r$, cohomology in degree $-n$ of its summand
is $H^{-n+r}(S\otL_D F^rN)$. The largest such $n$ is therefore
$\pd_D(F^rN)+r$. Cohomology commutes with this finite direct sum,
so the largest nonzero degree detected over $A$ is the maximum of
these values. The minimal-resolution formula over $A$ gives the result.
\end{proof}
This is the inflated-module case of
\cite[Corollary~4.11]{MinamotoYamaura}, with sides reversed.

\subsection{Tensor iterates and linear bounds}

Fix $R,e,\alpha$ over $k$ as in \Cref{prop:selection}, construct $B,P$
by \Cref{thm:realization}, and choose $[a_*,b_*]$, $C,W,O,T,D,X$
as in \Cref{sec:ordinary-simulation}. Put $A=D\ltimes X$.
For a finite-dimensional $R$-module $V$, regard $M(V)$ as supported
on the first factor of $D$. Put $E=B\otimes_k C$ and
$Q(V)=O\otimes_B M(V)$, a left $E$-module. Since $O_B$ is
projective, this ordinary tensor computes its derived tensor.
Moreover $Q(V)\ne0$ for $V\ne0$: the nonzero projective module
$O_B$ is a direct sum of vertex projectives, each of which evaluates
to $V$. The module $O$ is nonzero because otherwise $P=0$, whereas
\eqref{eq:realization} applied to $Y_2$ has nonzero right-hand side.

For $j\ge0$, induction from \eqref{eq:compression} and
\eqref{eq:realization} gives
\begin{align}
F^{2j}M(V)&\simeq\bigoplus_{q=0}^{j}
 M(H^{\circ j}V)[jb_*+3q]^{\oplus\binom jq},\label{eq:even-iterates}\\
F^{2j+1}M(V)&\simeq\bigoplus_{q=0}^{j}
 Q(H^{\circ j}V)[jb_*+3q]^{\oplus\binom jq}.\label{eq:odd-iterates}
\end{align}
Here the binomial coefficient counts direct summands.
Take any $Y_m$ satisfying \eqref{eq:extinction}, and inflate
$N_m=M(Y_m)$ along $A\to D\to B$. Then
\[
F^rN_m\ne0\ (0\le r<2m),\qquad F^{2m}N_m=0.
\]
For $0\le j<m$, Proposition~\ref{prop:exact-pd} therefore gives
\begin{equation}\label{eq:exact-final}
\begin{split}
\pd_A N_m=\max_{0\le j<m}\bigl\{
 &(b_*+5)j+\pd_B M(H^{\circ j}Y_m),\\
 &(b_*+5)j+1+\pd_E Q(H^{\circ j}Y_m)\bigr\}.
\end{split}
\end{equation}
For an even iterate, the shift and its bar degree add to
$jb_*+3q+2j$; for an odd iterate they add to $jb_*+3q+2j+1$.
Both are maximal at $q=j$, giving \eqref{eq:exact-final}.
For $j<m$, both modules in that formula are nonzero. Their projective
dimensions lie in $[0,2]$ and $[0,l+2]$, respectively. The odd term
at $j=m-1$ gives the lower bound. Since $b_*+5>0$, every term is at
most $(b_*+5)(m-1)+l+3$, giving
\begin{equation}\label{eq:linear-bounds}
(b_*+5)(m-1)+1\le\pd_A N_m\le(b_*+5)(m-1)+l+3.
\end{equation}

\section{Radicals and test modules}\label{sec:refinements}
\begin{proposition}\label{prop:radical}
The algebra $A$ above is split basic, $J_A^7=0$, and
$A/J_A\cong k^{3l+6}$. Its primitive vertex corners are $k$.
For every nonzero finite-dimensional $R$-module $Y$, the inflated
module $M(Y)$ has Loewy length three. For $R$-modules $Y,Z$,
\[
\Hom_A(M(Y),M(Z))\cong\Hom_R(Y,Z).
\]
\end{proposition}
\begin{proof}
The arrow ideals give $J_B^3=0$, $J_C^2=0$, and split semisimple
quotients. Thus
\[
J_E=J_B\otimes_k C+B\otimes_kJ_C,\qquad J_E^4=0.
\]
Indeed the displayed sum is nilpotent and its quotient is a product
of copies of $k$; any product of four terms contains three $J_B$
factors or two $J_C$ factors. Also $J_D=J_B\times J_E$ and
\[
J_A=J_D\oplus X,\qquad
J_A^n=J_D^n+\sum_{i+j=n-1}J_D^iXJ_D^j\quad(n\ge1).
\]
A word containing two $X$ factors vanishes because $XDX=0$.
The summand $O$ has left $E$-action and right $B$-action, so
$J_D^iOJ_D^j=0$ if $i\ge4$ or $j\ge3$. For $T$ the bounds are
reversed. Thus every term for $n=7$ vanishes.
There are $3+3(l+1)=3l+6$ vertices. Since $X$ joins distinct factors
of $D$, it contributes nothing to a primitive vertex corner;
each such corner of $A$ is $k$.

The action on $M(Y)$ factors through $B$, and the image of $J_A$
is $J_B$. The path $s's$ acts as identity from vertex $0$ to vertex
$2$. Therefore $J_A^2M(Y)\ne0$ and $J_A^3M(Y)=0$.
A $B$-module map $M(Y)\to M(Z)$ has three components. Commutation
with $s,s'$ makes them equal, and commutation with the arrows
$a_h$ makes their common value $R$-linear. Inflation preserves
these morphisms, proving the Hom identity.
\end{proof}

\subsection{Explicit complex representations}
For $k=\C$ use $p=2$. For $m\ge1$ set
$S_m=\F_2[t]/(t^m-1)$. Each element has a unique representative
$\sum_{r=0}^{m-1}a_rt^r$. Define $\lambda_m:S_m\to\F_2$ by
$\lambda_m(\sum a_rt^r)=a_{m-1}$. On the complex vector space
\[
Y_m^{\mathrm{exp}}=\bigoplus_{x\in S_m^3}\C v_x
\]
the vectors $v_x$ form a basis indexed by $S_m^3$.
Let $I$ be the $3\times3$ identity matrix. Write $e_i$ for the $i$th coordinate vector in $S_m^3$, $E_{ij}$
for the $3\times3$ matrix unit.
Define, for $i,j,\ell\in\{1,2,3\}$, $i\ne j$, and $r\in\Z$,
\begin{align*}
T_\ell v_x&=v_{x+(t-1)x_\ell e_\ell},&
U_i(r)v_x&=(-1)^{\lambda_m(t^rx_i)}v_x,\\
V_i(r)v_x&=v_{x+t^re_i},&
W_{ij}(r)v_x&=v_{(I+t^rE_{ij})x}.
\end{align*}
The exponent of $-1$ is read as $0$ or $1$.

\begin{proposition}\label{prop:explicit}
These operators define a simple $R$-module of dimension $2^{3m}$
satisfying \eqref{eq:extinction}.
\end{proposition}
\begin{proof}
Conjugation by $T_\ell$ changes the parameter of $U_i$ by
$-\delta_{\ell i}$, that of $V_i$ by $\delta_{\ell i}$, and that
of $W_{ij}$ by $\delta_{\ell i}-\delta_{\ell j}$.
Computing on $v_x$ gives
\[
[U_i(r),W_{ij}(s)]=U_j(r+s),\qquad
[W_{ij}(r),V_j(s)]=V_i(r+s),
\]
and $[U_i(r),V_i(s)]$ acts as
$(-1)^{\lambda_m(t^{r+s})}$. The remaining commutations follow
from their coordinate supports, and $U_i(r)^2=1$.
Thus all defining relations hold. The elements $z_0,\ldots,z_{m-2}$
act as $1$ and $z_{m-1}$ as $-1$, giving the required extinction.

The pairing $(a,b)\mapsto\lambda_m(ab)$ on $S_m$ is nondegenerate:
in the basis $1,t,\ldots,t^{m-1}$ its matrix has entry $1$ exactly
when the two exponents sum to $m-1$ modulo $m$, a permutation matrix.
For $a\in S_m^3$, products of the operators $U_i(r)$ therefore give
the diagonal operator
$U(a)v_y=(-1)^{\lambda_m(\sum_i a_iy_i)}v_y$.
For fixed $x\in S_m^3$, the operator
\[
2^{-3m}\sum_{a\in S_m^3}
(-1)^{\lambda_m(\sum_i a_ix_i)}U(a)
\]
acts on $v_y$ by the scalar
$2^{-3m}\sum_a(-1)^{\lambda_m(\sum_i a_i(x_i+y_i))}$.
If $y=x$, every term is $1$. If $y\ne x$, nondegeneracy supplies
$a_0$ with $\lambda_m(\sum_i(a_0)_i(x_i+y_i))=1$; translation
$a\mapsto a+a_0$ pairs terms of opposite signs. The operator is
therefore the projection onto $\C v_x$.
A nonzero vector in an invariant subspace has a nonzero coefficient
at some $v_x$, so this projection puts $v_x$ in the subspace.
The elements $1,t,\ldots,t^{m-1}$ span $S_m$ additively, and products
of the $V_i(r)$ translate $x$ to any element of $S_m^3$. The subspace
is consequently all of $Y_m^{\mathrm{exp}}$.
\end{proof}

\begin{proof}[Proof of \Cref{thm:main}]
Use the simple modules from \Cref{prop:selection}, or those from
\Cref{prop:explicit} over $\C$. Schur's lemma makes their endomorphism
rings division rings, and over $\C$ these rings are $\C$.
Different extinction times imply that the simple modules are
nonisomorphic, so their mutual Hom spaces vanish.
Proposition~\ref{prop:radical} transports these conclusions and gives
Loewy length three and all assertions about $A$. Each encoding has
three copies of its input space, giving $3\cdot2^{3m}$ in the explicit
case. Finally \eqref{eq:linear-bounds} proves the claimed bounds and
$\findim A=\infty$.
\end{proof}

\subsection{Size and other odd shifts}
With $d$ generators in the quadratic presentation of $R$ and relation
basis $\mathcal R$, counting vertices, arrows, and length-two paths gives
\[
\dim_k B=3+2(d+1)+(d+1)^2-|\mathcal R|.
\]
Since $\dim_k C=2l+1$ and $\dim_k W=1$, the construction gives
\[
\dim_k A=(2l+3)\dim_k B+\sum_n\dim_k P^n.
\]
The support and the terms of $P$ depend on the chosen lifts.

\begin{proposition}
For every odd integer $s\ge3$, the realization theorem remains valid
with $[3]$ replaced by $[s]$, with a complex $P$ chosen for that $s$.
\end{proposition}
\begin{proof}
In the notation of \Cref{lem:odd-double}, suppose an object
$S_{s-2}$ represents $U\oplus U[s-2]$, and $S_1$ represents
$U\oplus U[1]$. Take the map $S_{s-2}[1]\to S_1$ that is identity
on their common $U[1]$ and zero on the other summands. Fullness of
the embedding into the idempotent completion supplies the map in
the original category. Its cone represents $U\oplus U[s]$.
Starting with $S_1$ proves the assertion for every odd $s$.
The same finite diagram lifting and rectification apply.
The final truncation triangle splits because
$\Ext_B^{s+1}(M(HY),M(HY))=0$ for $s\ge3$.
\end{proof}

For $\Phi=P\otL_B-$ as above and $j\ge0$, let $\ker\Phi^j$ be the
full subcategory of objects $Z$ with $\Phi^j(Z)\simeq0$.
These are thick subcategories of $\D^b(B\text{-}\modu)$:
they are closed under shifts, cones, and direct summands because $\Phi$
is an exact additive functor. These kernels form a strictly increasing chain
$\ker\Phi^j\subsetneq\ker\Phi^{j+1}$ for $j\ge0$, where
$\Phi^0$ is the identity. Indeed $M(Y_{j+1})$ lies in the latter
kernel and not in the former, by the iteration formula.

\section{Left--right asymmetry}
The big finitistic dimension $\Findim A$ is defined as $\findim A$
with all left modules allowed in the supremum. Thus $\Findim A\ge\findim A$.
Injective modules generate when the smallest triangulated subcategory
of the unbounded derived category containing them and closed under
arbitrary coproducts is the whole category. Rickard's theorem
\cite[Theorem~4.3]{Rickard2019}, applied to the opposite algebra,
implies that injective left modules do not generate for the algebra
$A$ of Theorem~\ref{thm:main}.

The following consequence over $\C$ is retained from \cite{OpenAI198}.
\begin{corollary}
There is a finite-dimensional complex algebra $\Lambda$ such that
\[
\findim\Lambda=\Findim\Lambda=\infty,\qquad
\findim(\Lambda^{\mathrm{op}})=\Findim(\Lambda^{\mathrm{op}})=0.
\]
Injective left $\Lambda$-modules do not generate, whereas injective
right $\Lambda$-modules generate.
\end{corollary}
\begin{proof}
Take $k=\C$. The algebra $A^{\mathrm{op}}$ is basic by
Proposition~\ref{prop:radical}. Choose its primitive vertex
idempotents $e_1,\ldots,e_n$, and let
$S_i=e_iA^{\mathrm{op}}/e_iJ_{A^{\mathrm{op}}}$ be the corresponding
simple right modules. Set
\[
E_{\mathrm{aux}}=\prod_{i=1}^n\C[\epsilon_i]/(\epsilon_i^2),
\quad U_{\mathrm{aux}}=\bigoplus_{i=1}^n S_i,
\quad \widetilde\Lambda=
\begin{pmatrix}A^{\mathrm{op}}&0\\U_{\mathrm{aux}}&E_{\mathrm{aux}}\end{pmatrix}.
\]
The left $E_{\mathrm{aux}}$-action on $U_{\mathrm{aux}}$ factors
through $E_{\mathrm{aux}}/J_{E_{\mathrm{aux}}}=\C^n$, with its $i$th
factor acting on $S_i$. This is \cite[Construction~2.1]{Cummings2024}.
Set $\Lambda=\widetilde\Lambda^{\mathrm{op}}$. We give the module
calculations underlying \cite[Propositions~3.2--3.3 and 4.4]{Cummings2024}.

A projective right $A^{\mathrm{op}}$-module, placed on the old
vertices with zero component on the new vertices, is projective over
$\widetilde\Lambda$. Inflation preserves exactness and radical
differentials, so an inflated minimal projective resolution remains
minimal. Its length is unchanged. The finite-dimensional left
$A$-modules $N_m$ therefore give finite-dimensional right
$\widetilde\Lambda$-modules of the same unbounded finite projective
dimensions. Equivalently, $\findim\Lambda=\Findim\Lambda=\infty$.

For each $i$, let $f_i$ be the idempotent of the $i$th factor of
$E_{\mathrm{aux}}$. As a right $\widetilde\Lambda$-module,
$f_i\widetilde\Lambda$ has radical $S_i\oplus\C\epsilon_i$
and one-dimensional top supported at $f_i$. Its radical is
annihilated by $J_{\widetilde\Lambda}$. Let $L_i$ be the simple right
$\Lambda$-module at $f_i$, let $\widehat S_i$ be the dual of $S_i$
supported on the old vertices, and put
$I_i=\Hom_{\C}(f_i\widetilde\Lambda,\C)$, an injective right
$\Lambda$-module. Dualizing the radical sequence over $\C$ gives
the inclusion $u_i$ and quotient map $v_i$ in the exact sequence
\[
 0\longrightarrow L_i\xrightarrow{u_i}I_i
 \xrightarrow{v_i}\widehat S_i\oplus L_i\longrightarrow0.
\]
The modules $\widehat S_i,L_i$ exhaust the simple right
$\Lambda$-modules. In particular, each simple right module is a
quotient of an injective module.

This quotient property forces every right $\Lambda$-module of finite
projective dimension to be projective. Indeed, finite-dimensional
algebras have projective covers for arbitrary modules. If a module
had finite projective dimension $d>0$, its minimal resolution would
end in an inclusion $0\ne K\hookrightarrow P$ with $K$ projective
and $K\subseteq PJ_\Lambda$. Nilpotence of $J_\Lambda$ implies
$K/KJ_\Lambda\ne0$, so choose a simple quotient $q:K\to S$.
Choose an injective module mapping onto $S$. Projectivity of $K$
lifts $q$ to this injective module, and injectivity extends that lift
to $P$. The resulting map $P\to S$ restricts to $q$ on $K$ but
annihilates $PJ_\Lambda$, a contradiction. Thus
$\findim(\Lambda^{\mathrm{op}})=\Findim(\Lambda^{\mathrm{op}})=0$.

To check injective generation on the right, fix $i$ and let
$\operatorname{pr}_{L_i}$ and $\operatorname{pr}_{\widehat S_i}$
be the projections of $\widehat S_i\oplus L_i$. Set
$\delta_i=u_i\operatorname{pr}_{L_i}v_i:I_i\to I_i$.
The equality $v_iu_i=0$ implies
\[
 \delta_i^2=0,\qquad
 \operatorname{im}\delta_i=u_iL_i,\qquad
 \ker\delta_i=v_i^{-1}(\widehat S_i\oplus0).
\]
Let $K_i$ have term $I_i$ in degrees $n\le0$, zero in degrees
$n>0$, and differential $\delta_i$ in every negative degree.
Define a zero-differential complex $H_i$ by
\[
 H_i^n=\begin{cases}
 \widehat S_i,&n<0,\\
 \widehat S_i\oplus L_i,&n=0,\\
 0,&n>0.
 \end{cases}
\]
The maps $K_i^n\to H_i^n$ given by
$\operatorname{pr}_{\widehat S_i}v_i$ for $n<0$ and by $v_i$
for $n=0$ form a chain map, since $v_i\delta_i=0$.
In negative degrees its induced map is the isomorphism
$\ker\delta_i/u_iL_i\to\widehat S_i$; in degree zero it is
$I_i/u_iL_i\to\widehat S_i\oplus L_i$. Hence $K_i\simeq H_i$.

For $a\ge0$, let $K_{i,a}$ be the subcomplex of $K_i$ supported
in degrees $-a,\ldots,0$, with the inherited differentials.
Each is a bounded complex of injectives. The inclusions
$K_{i,a}\to K_{i,a+1}$ have colimit $K_i$. Let $\mathrm{shift}$
send summand $a$ of $\bigoplus_{a\ge0}K_{i,a}$ into summand $a+1$
by that inclusion. There is a degreewise exact sequence of complexes
\[
 0\longrightarrow\bigoplus_{a\ge0}K_{i,a}
 \xrightarrow{\,1-\mathrm{shift}\,}\bigoplus_{a\ge0}K_{i,a}
 \longrightarrow K_i\longrightarrow0,
\]
The associated triangle puts $K_i$ in the
triangulated subcategory generated by injectives and coproducts.
Such a subcategory is closed under direct summands: the same
telescope construction applied to a repeated idempotent on an
object gives its image. Since $H_i$ has zero differential, both
$\widehat S_i$ and $L_i$ are direct summands of $H_i$.
Thus all simple right modules belong to this subcategory. A
composition series then puts the regular module $\Lambda$ in it;
the regular module generates the unbounded derived category.
Injective right $\Lambda$-modules therefore generate.
Finally, Rickard's theorem and $\Findim\Lambda=\infty$ imply
that injective left $\Lambda$-modules do not generate.
\end{proof}

\begin{thebibliography}{10}

\bibitem{BalmerSchlichting2001}
Paul Balmer and Marco Schlichting.
\newblock Idempotent completion of triangulated categories.
\newblock {\em Journal of Algebra}, 236(2):819--834, 2001.
\newblock \href{https://doi.org/10.1006/jabr.2000.8529}{Publisher DOI};
  \href{https://webhomes.maths.ed.ac.uk/~v1ranick/papers/balmschl.pdf}{full
  text}.

\bibitem{Bass1960}
Hyman Bass.
\newblock Finitistic dimension and a homological generalization of semi-primary
  rings.
\newblock {\em Transactions of the American Mathematical Society},
  95(3):466--488, 1960.
\newblock \href{https://doi.org/10.1090/S0002-9947-1960-0157984-8}{Publisher
  DOI}.

\bibitem{Cummings2024}
Charley Cummings.
\newblock Left--right symmetry of finite finitistic dimension.
\newblock {\em Bulletin of the London Mathematical Society}, 56(2):624--633,
  2024.
\newblock \href{https://doi.org/10.1112/blms.12954}{Publisher DOI}.

\bibitem{Keller2000}
Bernhard Keller.
\newblock Bimodule complexes via strong homotopy actions.
\newblock {\em Algebras and Representation Theory}, 3(4):357--376, 2000.
\newblock \href{https://doi.org/10.1023/A:1009954126727}{Publisher DOI};
  \href{https://arxiv.org/abs/math/9910178v1}{arXiv:math/9910178v1}.

\bibitem{MinamotoYamaura}
Hiroyuki Minamoto and Kota Yamaura.
\newblock Homological dimension formulas for trivial extension algebras.
\newblock {\em Journal of Pure and Applied Algebra}, 224(8):106344, 2020.
\newblock 30 pp. \href{https://doi.org/10.1016/j.jpaa.2020.106344}{Publisher
  DOI}; \href{https://arxiv.org/abs/1710.01469v1}{arXiv:1710.01469v1}. All
  result numbers in citations to this paper refer to this arXiv version.

\bibitem{NRS2004}
Amnon Neeman, Andrew Ranicki, and Aidan Schofield.
\newblock Representations of algebras as universal localizations.
\newblock {\em Mathematical Proceedings of the Cambridge Philosophical
  Society}, 136(1):105--117, 2004.
\newblock \href{https://doi.org/10.1017/S030500410300700X}{Publisher DOI};
  \href{https://arxiv.org/abs/math/0205034v2}{arXiv:math/0205034v2}.

\bibitem{OpenAI198}
{OpenAI}.
\newblock An algebra of infinite little finitistic dimension.
\newblock
  \href{https://github.com/openai/math/tree/adc7f1241b42e322a6451854ab7e4b4c146bf78a/preprints/An-algebra-of-infinite-little-finitistic-dimension-September-23-2026}{OpenAI
  Math preprint}, 2026.
\newblock September 23, 2026. Source version: repository commit
  \texttt{adc7f1241b42e322a6451854ab7e4b4c146bf78a}.

\bibitem{Rickard2019}
Jeremy Rickard.
\newblock Unbounded derived categories and the finitistic dimension conjecture.
\newblock {\em Advances in Mathematics}, 354:106735, 2019.
\newblock \href{https://doi.org/10.1016/j.aim.2019.106735}{Publisher DOI};
  \href{https://arxiv.org/abs/1804.09801v1}{arXiv:1804.09801v1}.

\bibitem{Rego2022}
Yuri Santos~Rego.
\newblock On the finiteness length of some soluble linear groups.
\newblock {\em Canadian Journal of Mathematics}, 74(5):1209--1243, 2022.
\newblock \href{https://doi.org/10.4153/S0008414X21000213}{Publisher DOI};
  \href{https://arxiv.org/abs/1901.06704v3}{arXiv:1901.06704v3}.

\bibitem{StacksProject}
{The Stacks Project Authors}.
\newblock The {Stacks} project.
\newblock \url{https://stacks.math.columbia.edu}, 2026.
\newblock Tags \href{https://stacks.math.columbia.edu/tag/04VI}{04VI},
  \href{https://stacks.math.columbia.edu/tag/04VJ}{04VJ},
  \href{https://stacks.math.columbia.edu/tag/05RG}{05RG}, and
  \href{https://stacks.math.columbia.edu/tag/05RJ}{05RJ}; accessed 7 October 2026.

\bibitem{HuisgenZimmermann1995}
Birge Zimmermann~Huisgen.
\newblock The finitistic dimension conjectures---a tale of 3.5 decades.
\newblock In Alberto Facchini and Claudia Menini, editors, {\em Abelian Groups
  and Modules}, volume 343 of {\em Mathematics and Its Applications}, pages
  501--517. Kluwer Academic Publishers, Dordrecht, 1995.
\newblock \href{https://doi.org/10.1007/978-94-011-0443-2_41}{Publisher DOI};
  \href{https://arxiv.org/abs/1407.2383v1}{arXiv:1407.2383v1}.

\end{thebibliography}

\bigskip
\noindent Qiyue Tang\\
Zhili College, Tsinghua University\\
Beijing 100084, China\\
\textit{Email:} \href{mailto:tangqy24@mails.tsinghua.edu.cn}{\texttt{tangqy24@mails.tsinghua.edu.cn}}

\end{document}
